% Môn: Vật lý
% Lớp: 10
% Chương: Chương VI. Chuyển động tròn
% Bài: L10C6 Bài 31. Động học của chuyển động tròn đều
% Số câu: 164
% App đọc file này trực tiếp — sửa rồi Commit trên GitHub, bấm Đồng bộ GitHub .tex

% L10C6 Bài 31. Động học của chuyển động tròn đều

% ===== Câu 1 =====
% ID: L10C6B31-01-TN
% Mức: TH
\dangbt{Khái niệm cơ bản về chuyển động tròn đều}
\begin{ex}
% ID: L10C6B31-01-TN
% Mức: TH
Chuyển động nào dưới đây không phải là chuyển động tròn đều?
\choice
{Chuyển động của điểm đầu cánh quạt trần khi quạt đang quay.}
{\True Chuyển động của điểm đầu cánh quạt khi máy bay đang bay thẳng đều đối với người dưới đất.}
{Chuyển động của chiếc ống bương chứa nước trong cái cọn nước.}
{Chuyển động của con ngựa trong chiếc đu quay khi đang hoạt động ổn định.}
\loigiai{
Chuyển động tròn đều là chuyển động có quỹ đạo là đường tròn và tốc độ góc không đổi. Đối với chuyển động tròn đều, gia tốc hướng tâm luôn tồn tại và có độ lớn không đổi. Để xác định chuyển động nào không phải là chuyển động tròn đều, ta cần xem xét quỹ đạo và tốc độ góc của từng trường hợp.
 
 **A.** Sai — Chuyển động của điểm đầu cánh quạt trần khi quạt đang quay là chuyển động tròn đều vì quỹ đạo là đường tròn và tốc độ góc không đổi.
 **B.** Đúng — Chuyển động của điểm đầu cánh quạt khi máy bay đang bay thẳng đều đối với người dưới đất không phải là chuyển động tròn đều. Mặc dù quạt quay tròn đều so với máy bay, nhưng đối với người dưới đất, quỹ đạo của điểm đầu cánh quạt là một đường phức tạp, không phải đường tròn.
 **C.** Sai — Chuyển động của chiếc ống bương chứa nước trong cái cọn nước là chuyển động tròn đều vì quỹ đạo là đường tròn và tốc độ góc không đổi.
 **D.** Sai — Chuyển động của con ngựa trong chiếc đu quay khi đang hoạt động ổn định là chuyển động tròn đều vì quỹ đạo là đường tròn và tốc độ góc không đổi.
}
\end{ex}

% ===== Câu 2 =====
% ID: L10C6B31-02-TN
% Mức: NB
\dangbt{Tính chất chu kì trong chuyển động tròn đều}
\begin{ex}
% ID: L10C6B31-02-TN
% Mức: NB
Chu kì của kim phút là
\choice
{$1$ phút.}
{$360$ s.}
{\True $60$ phút.}
{$10$ phút.}
\loigiai{
**Phân tích đề bài:**
 
 Đề bài hỏi về chu kỳ của kim phút. Chu kỳ trong chuyển động tròn đều là thời gian để vật thực hiện được một vòng quay trọn vẹn.
 
 **Công thức liên quan:**
 
 Trong chuyển động tròn đều, chu kỳ ($T$) được định nghĩa là thời gian để hoàn thành một vòng quay.
 
 **Giải thích:**
 
 Kim phút của đồng hồ là một ví dụ về chuyển động tròn đều. Để kim phút quay hết một vòng tròn (tức là quay được 360 độ), nó cần di chuyển từ số này đến số kia trên mặt đồng hồ. Mỗi khoảng giữa hai số liên tiếp trên mặt đồng hồ tương ứng với 5 phút. Mặt đồng hồ có 12 số, do đó có 12 khoảng.
 
 Thời gian để kim phút quay hết một vòng là:
 $T = 12 \text{ khoảng} \times 5 \text{ phút/khoảng} = 60 \text{ phút}$ (Ngoài ra, chúng ta có thể biểu diễn chu kỳ theo đơn vị giây:) $T = 60 \text{ phút} \times 60 \text{ giây/phút} = 3600 \text{ giây}$ (**Kiểm tra các phương án:**) $60$ phút: **Đúng. Kim phút quay hết một vòng trong 60 phút.**
}
\end{ex}

% ===== Câu 3 =====
% ID: AIQ_L10C6_FULL_0001
% Mức: NB
\dangbt{Nhận biết chuyển động tròn và chuyển động tròn đều}
\begin{ex}
% ID: AIQ_L10C6_FULL_0001
% Mức: NB
Một vật đi được 1 vòng tròn bán kính $0,5$ m. Quãng đường vật đi được là bao nhiêu?
\choice
{\True $3,14\ \mathrm{m}$}
{$1,57\ \mathrm{m}$}
{$1\ \mathrm{m}$}
{Không đủ dữ kiện để có giá trị này}
\loigiai{
$s=n\cdot2\pi r=1\cdot2\pi\cdot0,5$. Suy ra kết quả $3,14\ \mathrm{m}$.
}
\end{ex}

% ===== Câu 4 =====
% ID: AIQ_L10C6_FULL_0002
% Mức: NB
\begin{ex}
% ID: AIQ_L10C6_FULL_0002
% Mức: NB
Một vật đi được 2 vòng tròn bán kính $0,6$ m. Quãng đường vật đi được là bao nhiêu?
\choice
{$3,77\ \mathrm{m}$}
{\True $7,54\ \mathrm{m}$}
{$2,4\ \mathrm{m}$}
{$3,77\ \mathrm{m}$}
\loigiai{
$s=n\cdot2\pi r=2\cdot2\pi\cdot0,6$. Suy ra kết quả $7,54\ \mathrm{m}$.
}
\end{ex}

% ===== Câu 5 =====
% ID: AIQ_L10C6_FULL_0003
% Mức: TH
\begin{ex}
% ID: AIQ_L10C6_FULL_0003
% Mức: TH
Một vật đi được 3 vòng tròn bán kính $0,7$ m. Quãng đường vật đi được là bao nhiêu?
\choice
{$6,6\ \mathrm{m}$}
{$4,2\ \mathrm{m}$}
{\True $13,19\ \mathrm{m}$}
{$4,4\ \mathrm{m}$}
\loigiai{
$s=n\cdot2\pi r=3\cdot2\pi\cdot0,7$. Suy ra kết quả $13,19\ \mathrm{m}$.
}
\end{ex}

% ===== Câu 6 =====
% ID: AIQ_L10C6_FULL_0004
% Mức: TH
\begin{ex}
% ID: AIQ_L10C6_FULL_0004
% Mức: TH
Một vật đi được 4 vòng tròn bán kính $0,8$ m. Quãng đường vật đi được là bao nhiêu?
\choice
{$10,05\ \mathrm{m}$}
{$6,4\ \mathrm{m}$}
{$5,03\ \mathrm{m}$}
{\True $20,11\ \mathrm{m}$}
\loigiai{
$s=n\cdot2\pi r=4\cdot2\pi\cdot0,8$. Suy ra kết quả $20,11\ \mathrm{m}$.
}
\end{ex}

% ===== Câu 7 =====
% ID: AIQ_L10C6_FULL_0005
% Mức: VD
\begin{ex}
% ID: AIQ_L10C6_FULL_0005
% Mức: VD
Một vật đi được 5 vòng tròn bán kính $0,9$ m. Quãng đường vật đi được là bao nhiêu?
\choice
{\True $28,27\ \mathrm{m}$}
{$14,14\ \mathrm{m}$}
{$9\ \mathrm{m}$}
{$5,65\ \mathrm{m}$}
\loigiai{
$s=n\cdot2\pi r=5\cdot2\pi\cdot0,9$. Suy ra kết quả $28,27\ \mathrm{m}$.
}
\end{ex}

% ===== Câu 8 =====
% ID: AIQ_L10C6_FULL_0006
% Mức: VD
\begin{ex}
% ID: AIQ_L10C6_FULL_0006
% Mức: VD
Một vật đi được 6 vòng tròn bán kính $0,5$ m. Quãng đường vật đi được là bao nhiêu?
\choice
{$9,42\ \mathrm{m}$}
{\True $18,85\ \mathrm{m}$}
{$6\ \mathrm{m}$}
{$3,14\ \mathrm{m}$}
\loigiai{
$s=n\cdot2\pi r=6\cdot2\pi\cdot0,5$. Suy ra kết quả $18,85\ \mathrm{m}$.
}
\end{ex}

% ===== Câu 9 =====
% ID: AIQ_L10C6_FULL_0007
% Mức: VD
\begin{ex}
% ID: AIQ_L10C6_FULL_0007
% Mức: VD
Một vật đi được 7 vòng tròn bán kính $0,6$ m. Quãng đường vật đi được là bao nhiêu?
\choice
{$13,19\ \mathrm{m}$}
{$8,4\ \mathrm{m}$}
{\True $26,39\ \mathrm{m}$}
{$3,77\ \mathrm{m}$}
\loigiai{
$s=n\cdot2\pi r=7\cdot2\pi\cdot0,6$. Suy ra kết quả $26,39\ \mathrm{m}$.
}
\end{ex}

% ===== Câu 10 =====
% ID: AIQ_L10C6_FULL_0008
% Mức: VDC
\begin{ex}
% ID: AIQ_L10C6_FULL_0008
% Mức: VDC
Một vật đi được 8 vòng tròn bán kính $0,7$ m. Quãng đường vật đi được là bao nhiêu?
\choice
{$17,59\ \mathrm{m}$}
{$11,2\ \mathrm{m}$}
{$4,4\ \mathrm{m}$}
{\True $35,19\ \mathrm{m}$}
\loigiai{
$s=n\cdot2\pi r=8\cdot2\pi\cdot0,7$. Suy ra kết quả $35,19\ \mathrm{m}$.
}
\end{ex}

% ===== Câu 11 =====
% ID: AIQ_L10C6_FULL_0009
% Mức: VDC
\begin{ex}
% ID: AIQ_L10C6_FULL_0009
% Mức: VDC
Một vật đi được 9 vòng tròn bán kính $0,8$ m. Quãng đường vật đi được là bao nhiêu?
\choice
{\True $45,24\ \mathrm{m}$}
{$22,62\ \mathrm{m}$}
{$14,4\ \mathrm{m}$}
{$5,03\ \mathrm{m}$}
\loigiai{
$s=n\cdot2\pi r=9\cdot2\pi\cdot0,8$. Suy ra kết quả $45,24\ \mathrm{m}$.
}
\end{ex}

% ===== Câu 12 =====
% ID: AIQ_L10C6_FULL_0010
% Mức: VDC
\begin{ex}
% ID: AIQ_L10C6_FULL_0010
% Mức: VDC
Một vật đi được 10 vòng tròn bán kính $0,9$ m. Quãng đường vật đi được là bao nhiêu?
\choice
{$28,27\ \mathrm{m}$}
{\True $56,55\ \mathrm{m}$}
{$18\ \mathrm{m}$}
{$5,65\ \mathrm{m}$}
\loigiai{
$s=n\cdot2\pi r=10\cdot2\pi\cdot0,9$. Suy ra kết quả $56,55\ \mathrm{m}$.
}
\end{ex}

% ===== Câu 13 =====
% ID: AIQ_L10C6_FULL_0011
% Mức: NB
\begin{ex}
% ID: AIQ_L10C6_FULL_0011
% Mức: NB
Xét các phát biểu thuộc dạng “Nhận biết chuyển động tròn và chuyển động tròn đều” ở tình huống số 1.
\choiceTF
{\True Chuyển động tròn có quỹ đạo là đường tròn.}
{\True Chuyển động tròn đều có tốc độ không đổi.}
{Vectơ vận tốc trong chuyển động tròn đều không đổi.}
{\True Vận tốc có phương tiếp tuyến quỹ đạo.}
\loigiai{
\begin{itemchoice}
\item (Đúng) Chuyển động tròn có quỹ đạo là đường tròn. (Vì): Phù hợp với định nghĩa hoặc hệ thức vật lí. b) Đúng: Phù hợp với định nghĩa hoặc hệ thức vật lí. c) Sai: Không phù hợp với định nghĩa hoặc hệ thức vật lí. d) Đúng: Phù hợp với định nghĩa hoặc hệ thức vật lí
\item (Đúng) Chuyển động tròn đều có tốc độ không đổi.
\item (Sai) Vectơ vận tốc trong chuyển động tròn đều không đổi.
\item (Đúng) Vận tốc có phương tiếp tuyến quỹ đạo.
\end{itemchoice}
}
\end{ex}

% ===== Câu 14 =====
% ID: AIQ_L10C6_FULL_0012
% Mức: TH
\begin{ex}
% ID: AIQ_L10C6_FULL_0012
% Mức: TH
Xét các phát biểu thuộc dạng “Nhận biết chuyển động tròn và chuyển động tròn đều” ở tình huống số 2.
\choiceTF
{\True Chuyển động tròn có quỹ đạo là đường tròn.}
{\True Chuyển động tròn đều có tốc độ không đổi.}
{Vectơ vận tốc trong chuyển động tròn đều không đổi.}
{\True Vận tốc có phương tiếp tuyến quỹ đạo.}
\loigiai{
\begin{itemchoice}
\item (Đúng) Chuyển động tròn có quỹ đạo là đường tròn. (Vì): Phù hợp với định nghĩa hoặc hệ thức vật lí. b) Đúng: Phù hợp với định nghĩa hoặc hệ thức vật lí. c) Sai: Không phù hợp với định nghĩa hoặc hệ thức vật lí. d) Đúng: Phù hợp với định nghĩa hoặc hệ thức vật lí
\item (Đúng) Chuyển động tròn đều có tốc độ không đổi.
\item (Sai) Vectơ vận tốc trong chuyển động tròn đều không đổi.
\item (Đúng) Vận tốc có phương tiếp tuyến quỹ đạo.
\end{itemchoice}
}
\end{ex}

% ===== Câu 15 =====
% ID: AIQ_L10C6_FULL_0013
% Mức: TH
\begin{ex}
% ID: AIQ_L10C6_FULL_0013
% Mức: TH
Xét các phát biểu thuộc dạng “Nhận biết chuyển động tròn và chuyển động tròn đều” ở tình huống số 3.
\choiceTF
{\True Chuyển động tròn có quỹ đạo là đường tròn.}
{\True Chuyển động tròn đều có tốc độ không đổi.}
{Vectơ vận tốc trong chuyển động tròn đều không đổi.}
{\True Vận tốc có phương tiếp tuyến quỹ đạo.}
\loigiai{
\begin{itemchoice}
\item (Đúng) Chuyển động tròn có quỹ đạo là đường tròn. (Vì): Phù hợp với định nghĩa hoặc hệ thức vật lí. b) Đúng: Phù hợp với định nghĩa hoặc hệ thức vật lí. c) Sai: Không phù hợp với định nghĩa hoặc hệ thức vật lí. d) Đúng: Phù hợp với định nghĩa hoặc hệ thức vật lí
\item (Đúng) Chuyển động tròn đều có tốc độ không đổi.
\item (Sai) Vectơ vận tốc trong chuyển động tròn đều không đổi.
\item (Đúng) Vận tốc có phương tiếp tuyến quỹ đạo.
\end{itemchoice}
}
\end{ex}

% ===== Câu 16 =====
% ID: AIQ_L10C6_FULL_0014
% Mức: VD
\begin{ex}
% ID: AIQ_L10C6_FULL_0014
% Mức: VD
Xét các phát biểu thuộc dạng “Nhận biết chuyển động tròn và chuyển động tròn đều” ở tình huống số 4.
\choiceTF
{\True Chuyển động tròn có quỹ đạo là đường tròn.}
{\True Chuyển động tròn đều có tốc độ không đổi.}
{Vectơ vận tốc trong chuyển động tròn đều không đổi.}
{\True Vận tốc có phương tiếp tuyến quỹ đạo.}
\loigiai{
\begin{itemchoice}
\item (Đúng) Chuyển động tròn có quỹ đạo là đường tròn. (Vì): Phù hợp với định nghĩa hoặc hệ thức vật lí. b) Đúng: Phù hợp với định nghĩa hoặc hệ thức vật lí. c) Sai: Không phù hợp với định nghĩa hoặc hệ thức vật lí. d) Đúng: Phù hợp với định nghĩa hoặc hệ thức vật lí
\item (Đúng) Chuyển động tròn đều có tốc độ không đổi.
\item (Sai) Vectơ vận tốc trong chuyển động tròn đều không đổi.
\item (Đúng) Vận tốc có phương tiếp tuyến quỹ đạo.
\end{itemchoice}
}
\end{ex}

% ===== Câu 17 =====
% ID: AIQ_L10C6_FULL_0015
% Mức: VDC
\begin{ex}
% ID: AIQ_L10C6_FULL_0015
% Mức: VDC
Xét các phát biểu thuộc dạng “Nhận biết chuyển động tròn và chuyển động tròn đều” ở tình huống số 5.
\choiceTF
{\True Chuyển động tròn có quỹ đạo là đường tròn.}
{\True Chuyển động tròn đều có tốc độ không đổi.}
{Vectơ vận tốc trong chuyển động tròn đều không đổi.}
{\True Vận tốc có phương tiếp tuyến quỹ đạo.}
\loigiai{
\begin{itemchoice}
\item (Đúng) Chuyển động tròn có quỹ đạo là đường tròn. (Vì): Phù hợp với định nghĩa hoặc hệ thức vật lí. b) Đúng: Phù hợp với định nghĩa hoặc hệ thức vật lí. c) Sai: Không phù hợp với định nghĩa hoặc hệ thức vật lí. d) Đúng: Phù hợp với định nghĩa hoặc hệ thức vật lí
\item (Đúng) Chuyển động tròn đều có tốc độ không đổi.
\item (Sai) Vectơ vận tốc trong chuyển động tròn đều không đổi.
\item (Đúng) Vận tốc có phương tiếp tuyến quỹ đạo.
\end{itemchoice}
}
\end{ex}

% ===== Câu 18 =====
% ID: AIQ_L10C6_FULL_0016
% Mức: TH
\begin{ex}
% ID: AIQ_L10C6_FULL_0016
% Mức: TH
Một vật đi được 11 vòng tròn bán kính $0,5$ m. Quãng đường vật đi được là bao nhiêu?
\shortans{34,56}
\loigiai{
$s=n\cdot2\pi r=11\cdot2\pi\cdot0,5$. Suy ra kết quả $34,56\ \mathrm{m}$. Đáp án trả lời ngắn không ghi đơn vị.
}
\end{ex}

% ===== Câu 19 =====
% ID: AIQ_L10C6_FULL_0017
% Mức: TH
\begin{ex}
% ID: AIQ_L10C6_FULL_0017
% Mức: TH
Một vật đi được 12 vòng tròn bán kính $0,6$ m. Quãng đường vật đi được là bao nhiêu?
\shortans{45,24}
\loigiai{
$s=n\cdot2\pi r=12\cdot2\pi\cdot0,6$. Suy ra kết quả $45,24\ \mathrm{m}$. Đáp án trả lời ngắn không ghi đơn vị.
}
\end{ex}

% ===== Câu 20 =====
% ID: AIQ_L10C6_FULL_0018
% Mức: VD
\begin{ex}
% ID: AIQ_L10C6_FULL_0018
% Mức: VD
Một vật đi được 13 vòng tròn bán kính $0,7$ m. Quãng đường vật đi được là bao nhiêu?
\shortans{57,18}
\loigiai{
$s=n\cdot2\pi r=13\cdot2\pi\cdot0,7$. Suy ra kết quả $57,18\ \mathrm{m}$. Đáp án trả lời ngắn không ghi đơn vị.
}
\end{ex}

% ===== Câu 21 =====
% ID: AIQ_L10C6_FULL_0019
% Mức: VD
\begin{ex}
% ID: AIQ_L10C6_FULL_0019
% Mức: VD
Một vật đi được 14 vòng tròn bán kính $0,8$ m. Quãng đường vật đi được là bao nhiêu?
\shortans{70,37}
\loigiai{
$s=n\cdot2\pi r=14\cdot2\pi\cdot0,8$. Suy ra kết quả $70,37\ \mathrm{m}$. Đáp án trả lời ngắn không ghi đơn vị.
}
\end{ex}

% ===== Câu 22 =====
% ID: AIQ_L10C6_FULL_0020
% Mức: VDC
\begin{ex}
% ID: AIQ_L10C6_FULL_0020
% Mức: VDC
Một vật đi được 15 vòng tròn bán kính $0,9$ m. Quãng đường vật đi được là bao nhiêu?
\shortans{84,82}
\loigiai{
$s=n\cdot2\pi r=15\cdot2\pi\cdot0,9$. Suy ra kết quả $84,82\ \mathrm{m}$. Đáp án trả lời ngắn không ghi đơn vị.
}
\end{ex}

% ===== Câu 23 =====
% ID: AIQ_L10C6_FULL_0021
% Mức: VDC
\begin{ex}
% ID: AIQ_L10C6_FULL_0021
% Mức: VDC
Một vật đi được 16 vòng tròn bán kính $0,5$ m. Quãng đường vật đi được là bao nhiêu?
\shortans{50,27}
\loigiai{
$s=n\cdot2\pi r=16\cdot2\pi\cdot0,5$. Suy ra kết quả $50,27\ \mathrm{m}$. Đáp án trả lời ngắn không ghi đơn vị.
}
\end{ex}

% ===== Câu 24 =====
% ID: AIQ_L10C6_FULL_0022
% Mức: TH
\begin{bt}
% ID: AIQ_L10C6_FULL_0022
% Mức: TH
Trình bày cách giải: Một vật đi được 17 vòng tròn bán kính $0,6$ m. Quãng đường vật đi được là bao nhiêu?
\loigiai{
$s=n\cdot2\pi r=17\cdot2\pi\cdot0,6$. Suy ra kết quả $64,09\ \mathrm{m}$.
}
\end{bt}

% ===== Câu 25 =====
% ID: AIQ_L10C6_FULL_0023
% Mức: TH
\begin{bt}
% ID: AIQ_L10C6_FULL_0023
% Mức: TH
Trình bày cách giải: Một vật đi được 18 vòng tròn bán kính $0,7$ m. Quãng đường vật đi được là bao nhiêu?
\loigiai{
$s=n\cdot2\pi r=18\cdot2\pi\cdot0,7$. Suy ra kết quả $79,17\ \mathrm{m}$.
}
\end{bt}

% ===== Câu 26 =====
% ID: AIQ_L10C6_FULL_0024
% Mức: VD
\begin{bt}
% ID: AIQ_L10C6_FULL_0024
% Mức: VD
Trình bày cách giải: Một vật đi được 19 vòng tròn bán kính $0,8$ m. Quãng đường vật đi được là bao nhiêu?
\loigiai{
$s=n\cdot2\pi r=19\cdot2\pi\cdot0,8$. Suy ra kết quả $95,5\ \mathrm{m}$.
}
\end{bt}

% ===== Câu 27 =====
% ID: AIQ_L10C6_FULL_0025
% Mức: VD
\begin{bt}
% ID: AIQ_L10C6_FULL_0025
% Mức: VD
Trình bày cách giải: Một vật đi được 20 vòng tròn bán kính $0,9$ m. Quãng đường vật đi được là bao nhiêu?
\loigiai{
$s=n\cdot2\pi r=20\cdot2\pi\cdot0,9$. Suy ra kết quả $113,1\ \mathrm{m}$.
}
\end{bt}

% ===== Câu 28 =====
% ID: AIQ_L10C6_FULL_0026
% Mức: VDC
\begin{bt}
% ID: AIQ_L10C6_FULL_0026
% Mức: VDC
Trình bày cách giải: Một vật đi được 21 vòng tròn bán kính $0,5$ m. Quãng đường vật đi được là bao nhiêu?
\loigiai{
$s=n\cdot2\pi r=21\cdot2\pi\cdot0,5$. Suy ra kết quả $65,97\ \mathrm{m}$.
}
\end{bt}

% ===== Câu 29 =====
% ID: AIQ_L10C6_FULL_0027
% Mức: VDC
\begin{bt}
% ID: AIQ_L10C6_FULL_0027
% Mức: VDC
Trình bày cách giải: Một vật đi được 22 vòng tròn bán kính $0,6$ m. Quãng đường vật đi được là bao nhiêu?
\loigiai{
$s=n\cdot2\pi r=22\cdot2\pi\cdot0,6$. Suy ra kết quả $82,94\ \mathrm{m}$.
}
\end{bt}

% ===== Câu 30 =====
% ID: AIQ_L10C6_FULL_0028
% Mức: NB
\dangbt{Tính tốc độ góc, chu kì và tần số}
\begin{ex}
% ID: AIQ_L10C6_FULL_0028
% Mức: NB
Vật chuyển động tròn đều có chu kì $T=2$ s. Tính tốc độ góc.
\choice
{\True $3,14\ \mathrm{rad/s}$}
{$12,57\ \mathrm{rad/s}$}
{$0,5\ \mathrm{rad/s}$}
{$0,32\ \mathrm{rad/s}$}
\loigiai{
$\omega=2\pi/T=2\pi/2$. Suy ra kết quả $3,14\ \mathrm{rad/s}$.
}
\end{ex}

% ===== Câu 31 =====
% ID: AIQ_L10C6_FULL_0029
% Mức: NB
\begin{ex}
% ID: AIQ_L10C6_FULL_0029
% Mức: NB
Vật chuyển động tròn đều có chu kì $T=3$ s. Tính tốc độ góc.
\choice
{$18,85\ \mathrm{rad/s}$}
{\True $2,09\ \mathrm{rad/s}$}
{$0,33\ \mathrm{rad/s}$}
{$0,48\ \mathrm{rad/s}$}
\loigiai{
$\omega=2\pi/T=2\pi/3$. Suy ra kết quả $2,09\ \mathrm{rad/s}$.
}
\end{ex}

% ===== Câu 32 =====
% ID: AIQ_L10C6_FULL_0030
% Mức: TH
\begin{ex}
% ID: AIQ_L10C6_FULL_0030
% Mức: TH
Vật chuyển động tròn đều có chu kì $T=4$ s. Tính tốc độ góc.
\choice
{$25,13\ \mathrm{rad/s}$}
{$0,25\ \mathrm{rad/s}$}
{\True $1,57\ \mathrm{rad/s}$}
{$0,64\ \mathrm{rad/s}$}
\loigiai{
$\omega=2\pi/T=2\pi/4$. Suy ra kết quả $1,57\ \mathrm{rad/s}$.
}
\end{ex}

% ===== Câu 33 =====
% ID: AIQ_L10C6_FULL_0031
% Mức: TH
\begin{ex}
% ID: AIQ_L10C6_FULL_0031
% Mức: TH
Vật chuyển động tròn đều có chu kì $T=5$ s. Tính tốc độ góc.
\choice
{$31,42\ \mathrm{rad/s}$}
{$0,2\ \mathrm{rad/s}$}
{$0,8\ \mathrm{rad/s}$}
{\True $1,26\ \mathrm{rad/s}$}
\loigiai{
$\omega=2\pi/T=2\pi/5$. Suy ra kết quả $1,26\ \mathrm{rad/s}$.
}
\end{ex}

% ===== Câu 34 =====
% ID: AIQ_L10C6_FULL_0032
% Mức: VD
\begin{ex}
% ID: AIQ_L10C6_FULL_0032
% Mức: VD
Vật chuyển động tròn đều có chu kì $T=6$ s. Tính tốc độ góc.
\choice
{\True $1,05\ \mathrm{rad/s}$}
{$37,7\ \mathrm{rad/s}$}
{$0,17\ \mathrm{rad/s}$}
{$0,95\ \mathrm{rad/s}$}
\loigiai{
$\omega=2\pi/T=2\pi/6$. Suy ra kết quả $1,05\ \mathrm{rad/s}$.
}
\end{ex}

% ===== Câu 35 =====
% ID: AIQ_L10C6_FULL_0033
% Mức: VD
\begin{ex}
% ID: AIQ_L10C6_FULL_0033
% Mức: VD
Vật chuyển động tròn đều có chu kì $T=2$ s. Tính tốc độ góc.
\choice
{$12,57\ \mathrm{rad/s}$}
{\True $3,14\ \mathrm{rad/s}$}
{$0,5\ \mathrm{rad/s}$}
{$0,32\ \mathrm{rad/s}$}
\loigiai{
$\omega=2\pi/T=2\pi/2$. Suy ra kết quả $3,14\ \mathrm{rad/s}$.
}
\end{ex}

% ===== Câu 36 =====
% ID: AIQ_L10C6_FULL_0034
% Mức: VD
\begin{ex}
% ID: AIQ_L10C6_FULL_0034
% Mức: VD
Vật chuyển động tròn đều có chu kì $T=3$ s. Tính tốc độ góc.
\choice
{$18,85\ \mathrm{rad/s}$}
{$0,33\ \mathrm{rad/s}$}
{\True $2,09\ \mathrm{rad/s}$}
{$0,48\ \mathrm{rad/s}$}
\loigiai{
$\omega=2\pi/T=2\pi/3$. Suy ra kết quả $2,09\ \mathrm{rad/s}$.
}
\end{ex}

% ===== Câu 37 =====
% ID: AIQ_L10C6_FULL_0035
% Mức: VDC
\begin{ex}
% ID: AIQ_L10C6_FULL_0035
% Mức: VDC
Vật chuyển động tròn đều có chu kì $T=4$ s. Tính tốc độ góc.
\choice
{$25,13\ \mathrm{rad/s}$}
{$0,25\ \mathrm{rad/s}$}
{$0,64\ \mathrm{rad/s}$}
{\True $1,57\ \mathrm{rad/s}$}
\loigiai{
$\omega=2\pi/T=2\pi/4$. Suy ra kết quả $1,57\ \mathrm{rad/s}$.
}
\end{ex}

% ===== Câu 38 =====
% ID: AIQ_L10C6_FULL_0036
% Mức: VDC
\begin{ex}
% ID: AIQ_L10C6_FULL_0036
% Mức: VDC
Vật chuyển động tròn đều có chu kì $T=5$ s. Tính tốc độ góc.
\choice
{\True $1,26\ \mathrm{rad/s}$}
{$31,42\ \mathrm{rad/s}$}
{$0,2\ \mathrm{rad/s}$}
{$0,8\ \mathrm{rad/s}$}
\loigiai{
$\omega=2\pi/T=2\pi/5$. Suy ra kết quả $1,26\ \mathrm{rad/s}$.
}
\end{ex}

% ===== Câu 39 =====
% ID: AIQ_L10C6_FULL_0037
% Mức: VDC
\begin{ex}
% ID: AIQ_L10C6_FULL_0037
% Mức: VDC
Vật chuyển động tròn đều có chu kì $T=6$ s. Tính tốc độ góc.
\choice
{$37,7\ \mathrm{rad/s}$}
{\True $1,05\ \mathrm{rad/s}$}
{$0,17\ \mathrm{rad/s}$}
{$0,95\ \mathrm{rad/s}$}
\loigiai{
$\omega=2\pi/T=2\pi/6$. Suy ra kết quả $1,05\ \mathrm{rad/s}$.
}
\end{ex}

% ===== Câu 40 =====
% ID: AIQ_L10C6_FULL_0038
% Mức: NB
\begin{ex}
% ID: AIQ_L10C6_FULL_0038
% Mức: NB
Xét các phát biểu thuộc dạng “Tính tốc độ góc, chu kì và tần số” ở tình huống số 1.
\choiceTF
{\True $f=1/T$.}
{\True $\omega=2\pi/T$.}
{Đơn vị tốc độ góc là m/s.}
{\True Chu kì là thời gian đi một vòng.}
\loigiai{
\begin{itemchoice}
\item (Đúng) $f=1/T$. (Vì): Phù hợp với định nghĩa hoặc hệ thức vật lí. b) Đúng: Phù hợp với định nghĩa hoặc hệ thức vật lí. c) Sai: Không phù hợp với định nghĩa hoặc hệ thức vật lí. d) Đúng: Phù hợp với định nghĩa hoặc hệ thức vật lí
\item (Đúng) $\omega=2\pi/T$.
\item (Sai) Đơn vị tốc độ góc là m/s.
\item (Đúng) Chu kì là thời gian đi một vòng.
\end{itemchoice}
}
\end{ex}

% ===== Câu 41 =====
% ID: AIQ_L10C6_FULL_0039
% Mức: TH
\begin{ex}
% ID: AIQ_L10C6_FULL_0039
% Mức: TH
Xét các phát biểu thuộc dạng “Tính tốc độ góc, chu kì và tần số” ở tình huống số 2.
\choiceTF
{\True $f=1/T$.}
{\True $\omega=2\pi/T$.}
{Đơn vị tốc độ góc là m/s.}
{\True Chu kì là thời gian đi một vòng.}
\loigiai{
\begin{itemchoice}
\item (Đúng) $f=1/T$. (Vì): Phù hợp với định nghĩa hoặc hệ thức vật lí. b) Đúng: Phù hợp với định nghĩa hoặc hệ thức vật lí. c) Sai: Không phù hợp với định nghĩa hoặc hệ thức vật lí. d) Đúng: Phù hợp với định nghĩa hoặc hệ thức vật lí
\item (Đúng) $\omega=2\pi/T$.
\item (Sai) Đơn vị tốc độ góc là m/s.
\item (Đúng) Chu kì là thời gian đi một vòng.
\end{itemchoice}
}
\end{ex}

% ===== Câu 42 =====
% ID: AIQ_L10C6_FULL_0040
% Mức: TH
\begin{ex}
% ID: AIQ_L10C6_FULL_0040
% Mức: TH
Xét các phát biểu thuộc dạng “Tính tốc độ góc, chu kì và tần số” ở tình huống số 3.
\choiceTF
{\True $f=1/T$.}
{\True $\omega=2\pi/T$.}
{Đơn vị tốc độ góc là m/s.}
{\True Chu kì là thời gian đi một vòng.}
\loigiai{
\begin{itemchoice}
\item (Đúng) $f=1/T$. (Vì): Phù hợp với định nghĩa hoặc hệ thức vật lí. b) Đúng: Phù hợp với định nghĩa hoặc hệ thức vật lí. c) Sai: Không phù hợp với định nghĩa hoặc hệ thức vật lí. d) Đúng: Phù hợp với định nghĩa hoặc hệ thức vật lí
\item (Đúng) $\omega=2\pi/T$.
\item (Sai) Đơn vị tốc độ góc là m/s.
\item (Đúng) Chu kì là thời gian đi một vòng.
\end{itemchoice}
}
\end{ex}

% ===== Câu 43 =====
% ID: AIQ_L10C6_FULL_0041
% Mức: VD
\begin{ex}
% ID: AIQ_L10C6_FULL_0041
% Mức: VD
Xét các phát biểu thuộc dạng “Tính tốc độ góc, chu kì và tần số” ở tình huống số 4.
\choiceTF
{\True $f=1/T$.}
{\True $\omega=2\pi/T$.}
{Đơn vị tốc độ góc là m/s.}
{\True Chu kì là thời gian đi một vòng.}
\loigiai{
\begin{itemchoice}
\item (Đúng) $f=1/T$. (Vì): Phù hợp với định nghĩa hoặc hệ thức vật lí. b) Đúng: Phù hợp với định nghĩa hoặc hệ thức vật lí. c) Sai: Không phù hợp với định nghĩa hoặc hệ thức vật lí. d) Đúng: Phù hợp với định nghĩa hoặc hệ thức vật lí
\item (Đúng) $\omega=2\pi/T$.
\item (Sai) Đơn vị tốc độ góc là m/s.
\item (Đúng) Chu kì là thời gian đi một vòng.
\end{itemchoice}
}
\end{ex}

% ===== Câu 44 =====
% ID: AIQ_L10C6_FULL_0042
% Mức: VDC
\begin{ex}
% ID: AIQ_L10C6_FULL_0042
% Mức: VDC
Xét các phát biểu thuộc dạng “Tính tốc độ góc, chu kì và tần số” ở tình huống số 5.
\choiceTF
{\True $f=1/T$.}
{\True $\omega=2\pi/T$.}
{Đơn vị tốc độ góc là m/s.}
{\True Chu kì là thời gian đi một vòng.}
\loigiai{
\begin{itemchoice}
\item (Đúng) $f=1/T$. (Vì): Phù hợp với định nghĩa hoặc hệ thức vật lí. b) Đúng: Phù hợp với định nghĩa hoặc hệ thức vật lí. c) Sai: Không phù hợp với định nghĩa hoặc hệ thức vật lí. d) Đúng: Phù hợp với định nghĩa hoặc hệ thức vật lí
\item (Đúng) $\omega=2\pi/T$.
\item (Sai) Đơn vị tốc độ góc là m/s.
\item (Đúng) Chu kì là thời gian đi một vòng.
\end{itemchoice}
}
\end{ex}

% ===== Câu 45 =====
% ID: AIQ_L10C6_FULL_0043
% Mức: TH
\begin{ex}
% ID: AIQ_L10C6_FULL_0043
% Mức: TH
Vật chuyển động tròn đều có chu kì $T=2$ s. Tính tốc độ góc.
\shortans{3,14}
\loigiai{
$\omega=2\pi/T=2\pi/2$. Suy ra kết quả $3,14\ \mathrm{rad/s}$. Đáp án trả lời ngắn không ghi đơn vị.
}
\end{ex}

% ===== Câu 46 =====
% ID: AIQ_L10C6_FULL_0044
% Mức: TH
\begin{ex}
% ID: AIQ_L10C6_FULL_0044
% Mức: TH
Vật chuyển động tròn đều có chu kì $T=3$ s. Tính tốc độ góc.
\shortans{2,09}
\loigiai{
$\omega=2\pi/T=2\pi/3$. Suy ra kết quả $2,09\ \mathrm{rad/s}$. Đáp án trả lời ngắn không ghi đơn vị.
}
\end{ex}

% ===== Câu 47 =====
% ID: AIQ_L10C6_FULL_0045
% Mức: VD
\begin{ex}
% ID: AIQ_L10C6_FULL_0045
% Mức: VD
Vật chuyển động tròn đều có chu kì $T=4$ s. Tính tốc độ góc.
\shortans{1,57}
\loigiai{
$\omega=2\pi/T=2\pi/4$. Suy ra kết quả $1,57\ \mathrm{rad/s}$. Đáp án trả lời ngắn không ghi đơn vị.
}
\end{ex}

% ===== Câu 48 =====
% ID: AIQ_L10C6_FULL_0046
% Mức: VD
\begin{ex}
% ID: AIQ_L10C6_FULL_0046
% Mức: VD
Vật chuyển động tròn đều có chu kì $T=5$ s. Tính tốc độ góc.
\shortans{1,26}
\loigiai{
$\omega=2\pi/T=2\pi/5$. Suy ra kết quả $1,26\ \mathrm{rad/s}$. Đáp án trả lời ngắn không ghi đơn vị.
}
\end{ex}

% ===== Câu 49 =====
% ID: AIQ_L10C6_FULL_0047
% Mức: VDC
\begin{ex}
% ID: AIQ_L10C6_FULL_0047
% Mức: VDC
Vật chuyển động tròn đều có chu kì $T=6$ s. Tính tốc độ góc.
\shortans{1,05}
\loigiai{
$\omega=2\pi/T=2\pi/6$. Suy ra kết quả $1,05\ \mathrm{rad/s}$. Đáp án trả lời ngắn không ghi đơn vị.
}
\end{ex}

% ===== Câu 50 =====
% ID: AIQ_L10C6_FULL_0048
% Mức: VDC
\begin{ex}
% ID: AIQ_L10C6_FULL_0048
% Mức: VDC
Vật chuyển động tròn đều có chu kì $T=2$ s. Tính tốc độ góc.
\shortans{3,14}
\loigiai{
$\omega=2\pi/T=2\pi/2$. Suy ra kết quả $3,14\ \mathrm{rad/s}$. Đáp án trả lời ngắn không ghi đơn vị.
}
\end{ex}

% ===== Câu 51 =====
% ID: AIQ_L10C6_FULL_0049
% Mức: TH
\begin{ex}
% ID: AIQ_L10C6_FULL_0049
% Mức: TH
Trình bày cách giải: Vật chuyển động tròn đều có chu kì $T=3$ s. Tính tốc độ góc.
\choiceTF

\loigiai{
$\omega=2\pi/T=2\pi/3$. Suy ra kết quả $2,09\ \mathrm{rad/s}$.
}
\end{ex}

% ===== Câu 52 =====
% ID: AIQ_L10C6_FULL_0050
% Mức: TH
\begin{ex}
% ID: AIQ_L10C6_FULL_0050
% Mức: TH
Trình bày cách giải: Vật chuyển động tròn đều có chu kì $T=4$ s. Tính tốc độ góc.
\choiceTF

\loigiai{
$\omega=2\pi/T=2\pi/4$. Suy ra kết quả $1,57\ \mathrm{rad/s}$.
}
\end{ex}

% ===== Câu 53 =====
% ID: AIQ_L10C6_FULL_0051
% Mức: VD
\begin{ex}
% ID: AIQ_L10C6_FULL_0051
% Mức: VD
Trình bày cách giải: Vật chuyển động tròn đều có chu kì $T=5$ s. Tính tốc độ góc.
\choiceTF

\loigiai{
$\omega=2\pi/T=2\pi/5$. Suy ra kết quả $1,26\ \mathrm{rad/s}$.
}
\end{ex}

% ===== Câu 54 =====
% ID: AIQ_L10C6_FULL_0052
% Mức: VD
\begin{ex}
% ID: AIQ_L10C6_FULL_0052
% Mức: VD
Trình bày cách giải: Vật chuyển động tròn đều có chu kì $T=6$ s. Tính tốc độ góc.
\choiceTF

\loigiai{
$\omega=2\pi/T=2\pi/6$. Suy ra kết quả $1,05\ \mathrm{rad/s}$.
}
\end{ex}

% ===== Câu 55 =====
% ID: AIQ_L10C6_FULL_0053
% Mức: VDC
\begin{ex}
% ID: AIQ_L10C6_FULL_0053
% Mức: VDC
Trình bày cách giải: Vật chuyển động tròn đều có chu kì $T=2$ s. Tính tốc độ góc.
\choiceTF

\loigiai{
$\omega=2\pi/T=2\pi/2$. Suy ra kết quả $3,14\ \mathrm{rad/s}$.
}
\end{ex}

% ===== Câu 56 =====
% ID: AIQ_L10C6_FULL_0054
% Mức: VDC
\begin{ex}
% ID: AIQ_L10C6_FULL_0054
% Mức: VDC
Trình bày cách giải: Vật chuyển động tròn đều có chu kì $T=3$ s. Tính tốc độ góc.
\choiceTF

\loigiai{
$\omega=2\pi/T=2\pi/3$. Suy ra kết quả $2,09\ \mathrm{rad/s}$.
}
\end{ex}

% ===== Câu 57 =====
% ID: AIQ_L10C6_FULL_0055
% Mức: NB
\dangbt{Liên hệ tốc độ dài với tốc độ góc}
\begin{ex}
% ID: AIQ_L10C6_FULL_0055
% Mức: NB
Một điểm quay với tốc độ góc $2$ rad/s trên quỹ đạo bán kính $0,2$ m. Tính tốc độ dài.
\choice
{\True $0,4\ \mathrm{m/s}$}
{$10\ \mathrm{m/s}$}
{$0,1\ \mathrm{m/s}$}
{$2,2\ \mathrm{m/s}$}
\loigiai{
$v=\omega r=2\cdot0,2$. Suy ra kết quả $0,4\ \mathrm{m/s}$.
}
\end{ex}

% ===== Câu 58 =====
% ID: AIQ_L10C6_FULL_0056
% Mức: NB
\begin{ex}
% ID: AIQ_L10C6_FULL_0056
% Mức: NB
Một điểm quay với tốc độ góc $3$ rad/s trên quỹ đạo bán kính $0,3$ m. Tính tốc độ dài.
\choice
{$10\ \mathrm{m/s}$}
{\True $0,9\ \mathrm{m/s}$}
{$0,1\ \mathrm{m/s}$}
{$3,3\ \mathrm{m/s}$}
\loigiai{
$v=\omega r=3\cdot0,3$. Suy ra kết quả $0,9\ \mathrm{m/s}$.
}
\end{ex}

% ===== Câu 59 =====
% ID: AIQ_L10C6_FULL_0057
% Mức: TH
\begin{ex}
% ID: AIQ_L10C6_FULL_0057
% Mức: TH
Một điểm quay với tốc độ góc $4$ rad/s trên quỹ đạo bán kính $0,4$ m. Tính tốc độ dài.
\choice
{$10\ \mathrm{m/s}$}
{$0,1\ \mathrm{m/s}$}
{\True $1,6\ \mathrm{m/s}$}
{$4,4\ \mathrm{m/s}$}
\loigiai{
$v=\omega r=4\cdot0,4$. Suy ra kết quả $1,6\ \mathrm{m/s}$.
}
\end{ex}

% ===== Câu 60 =====
% ID: AIQ_L10C6_FULL_0058
% Mức: TH
\begin{ex}
% ID: AIQ_L10C6_FULL_0058
% Mức: TH
Một điểm quay với tốc độ góc $5$ rad/s trên quỹ đạo bán kính $0,5$ m. Tính tốc độ dài.
\choice
{$10\ \mathrm{m/s}$}
{$0,1\ \mathrm{m/s}$}
{$5,5\ \mathrm{m/s}$}
{\True $2,5\ \mathrm{m/s}$}
\loigiai{
$v=\omega r=5\cdot0,5$. Suy ra kết quả $2,5\ \mathrm{m/s}$.
}
\end{ex}

% ===== Câu 61 =====
% ID: AIQ_L10C6_FULL_0059
% Mức: VD
\begin{ex}
% ID: AIQ_L10C6_FULL_0059
% Mức: VD
Một điểm quay với tốc độ góc $6$ rad/s trên quỹ đạo bán kính $0,6$ m. Tính tốc độ dài.
\choice
{\True $3,6\ \mathrm{m/s}$}
{$10\ \mathrm{m/s}$}
{$0,1\ \mathrm{m/s}$}
{$6,6\ \mathrm{m/s}$}
\loigiai{
$v=\omega r=6\cdot0,6$. Suy ra kết quả $3,6\ \mathrm{m/s}$.
}
\end{ex}

% ===== Câu 62 =====
% ID: AIQ_L10C6_FULL_0060
% Mức: VD
\begin{ex}
% ID: AIQ_L10C6_FULL_0060
% Mức: VD
Một điểm quay với tốc độ góc $7$ rad/s trên quỹ đạo bán kính $0,2$ m. Tính tốc độ dài.
\choice
{$35\ \mathrm{m/s}$}
{\True $1,4\ \mathrm{m/s}$}
{$0,03\ \mathrm{m/s}$}
{$7,2\ \mathrm{m/s}$}
\loigiai{
$v=\omega r=7\cdot0,2$. Suy ra kết quả $1,4\ \mathrm{m/s}$.
}
\end{ex}

% ===== Câu 63 =====
% ID: AIQ_L10C6_FULL_0061
% Mức: VD
\begin{ex}
% ID: AIQ_L10C6_FULL_0061
% Mức: VD
Một điểm quay với tốc độ góc $2$ rad/s trên quỹ đạo bán kính $0,3$ m. Tính tốc độ dài.
\choice
{$6,67\ \mathrm{m/s}$}
{$0,15\ \mathrm{m/s}$}
{\True $0,6\ \mathrm{m/s}$}
{$2,3\ \mathrm{m/s}$}
\loigiai{
$v=\omega r=2\cdot0,3$. Suy ra kết quả $0,6\ \mathrm{m/s}$.
}
\end{ex}

% ===== Câu 64 =====
% ID: AIQ_L10C6_FULL_0062
% Mức: VDC
\begin{ex}
% ID: AIQ_L10C6_FULL_0062
% Mức: VDC
Một điểm quay với tốc độ góc $3$ rad/s trên quỹ đạo bán kính $0,4$ m. Tính tốc độ dài.
\choice
{$7,5\ \mathrm{m/s}$}
{$0,13\ \mathrm{m/s}$}
{$3,4\ \mathrm{m/s}$}
{\True $1,2\ \mathrm{m/s}$}
\loigiai{
$v=\omega r=3\cdot0,4$. Suy ra kết quả $1,2\ \mathrm{m/s}$.
}
\end{ex}

% ===== Câu 65 =====
% ID: AIQ_L10C6_FULL_0063
% Mức: VDC
\begin{ex}
% ID: AIQ_L10C6_FULL_0063
% Mức: VDC
Một điểm quay với tốc độ góc $4$ rad/s trên quỹ đạo bán kính $0,5$ m. Tính tốc độ dài.
\choice
{\True $2\ \mathrm{m/s}$}
{$8\ \mathrm{m/s}$}
{$0,13\ \mathrm{m/s}$}
{$4,5\ \mathrm{m/s}$}
\loigiai{
$v=\omega r=4\cdot0,5$. Suy ra kết quả $2\ \mathrm{m/s}$.
}
\end{ex}

% ===== Câu 66 =====
% ID: AIQ_L10C6_FULL_0064
% Mức: VDC
\begin{ex}
% ID: AIQ_L10C6_FULL_0064
% Mức: VDC
Một điểm quay với tốc độ góc $5$ rad/s trên quỹ đạo bán kính $0,6$ m. Tính tốc độ dài.
\choice
{$8,33\ \mathrm{m/s}$}
{\True $3\ \mathrm{m/s}$}
{$0,12\ \mathrm{m/s}$}
{$5,6\ \mathrm{m/s}$}
\loigiai{
$v=\omega r=5\cdot0,6$. Suy ra kết quả $3\ \mathrm{m/s}$.
}
\end{ex}

% ===== Câu 67 =====
% ID: AIQ_L10C6_FULL_0065
% Mức: NB
\begin{ex}
% ID: AIQ_L10C6_FULL_0065
% Mức: NB
Xét các phát biểu thuộc dạng “Liên hệ tốc độ dài với tốc độ góc” ở tình huống số 1.
\choiceTF
{\True $v=\omega r$.}
{\True Cùng tốc độ góc, điểm xa tâm có tốc độ dài lớn hơn.}
{Tốc độ dài có đơn vị rad/s.}
{\True Các điểm trên vật rắn quay có cùng tốc độ góc.}
\loigiai{
\begin{itemchoice}
\item (Đúng) $v=\omega r$. (Vì): Phù hợp với định nghĩa hoặc hệ thức vật lí. b) Đúng: Phù hợp với định nghĩa hoặc hệ thức vật lí. c) Sai: Không phù hợp với định nghĩa hoặc hệ thức vật lí. d) Đúng: Phù hợp với định nghĩa hoặc hệ thức vật lí
\item (Đúng) Cùng tốc độ góc, điểm xa tâm có tốc độ dài lớn hơn.
\item (Sai) Tốc độ dài có đơn vị rad/s.
\item (Đúng) Các điểm trên vật rắn quay có cùng tốc độ góc.
\end{itemchoice}
}
\end{ex}

% ===== Câu 68 =====
% ID: AIQ_L10C6_FULL_0066
% Mức: TH
\begin{ex}
% ID: AIQ_L10C6_FULL_0066
% Mức: TH
Xét các phát biểu thuộc dạng “Liên hệ tốc độ dài với tốc độ góc” ở tình huống số 2.
\choiceTF
{\True $v=\omega r$.}
{\True Cùng tốc độ góc, điểm xa tâm có tốc độ dài lớn hơn.}
{Tốc độ dài có đơn vị rad/s.}
{\True Các điểm trên vật rắn quay có cùng tốc độ góc.}
\loigiai{
\begin{itemchoice}
\item (Đúng) $v=\omega r$. (Vì): Phù hợp với định nghĩa hoặc hệ thức vật lí. b) Đúng: Phù hợp với định nghĩa hoặc hệ thức vật lí. c) Sai: Không phù hợp với định nghĩa hoặc hệ thức vật lí. d) Đúng: Phù hợp với định nghĩa hoặc hệ thức vật lí
\item (Đúng) Cùng tốc độ góc, điểm xa tâm có tốc độ dài lớn hơn.
\item (Sai) Tốc độ dài có đơn vị rad/s.
\item (Đúng) Các điểm trên vật rắn quay có cùng tốc độ góc.
\end{itemchoice}
}
\end{ex}

% ===== Câu 69 =====
% ID: AIQ_L10C6_FULL_0067
% Mức: TH
\begin{ex}
% ID: AIQ_L10C6_FULL_0067
% Mức: TH
Xét các phát biểu thuộc dạng “Liên hệ tốc độ dài với tốc độ góc” ở tình huống số 3.
\choiceTF
{\True $v=\omega r$.}
{\True Cùng tốc độ góc, điểm xa tâm có tốc độ dài lớn hơn.}
{Tốc độ dài có đơn vị rad/s.}
{\True Các điểm trên vật rắn quay có cùng tốc độ góc.}
\loigiai{
\begin{itemchoice}
\item (Đúng) $v=\omega r$. (Vì): Phù hợp với định nghĩa hoặc hệ thức vật lí. b) Đúng: Phù hợp với định nghĩa hoặc hệ thức vật lí. c) Sai: Không phù hợp với định nghĩa hoặc hệ thức vật lí. d) Đúng: Phù hợp với định nghĩa hoặc hệ thức vật lí
\item (Đúng) Cùng tốc độ góc, điểm xa tâm có tốc độ dài lớn hơn.
\item (Sai) Tốc độ dài có đơn vị rad/s.
\item (Đúng) Các điểm trên vật rắn quay có cùng tốc độ góc.
\end{itemchoice}
}
\end{ex}

% ===== Câu 70 =====
% ID: AIQ_L10C6_FULL_0068
% Mức: VD
\begin{ex}
% ID: AIQ_L10C6_FULL_0068
% Mức: VD
Xét các phát biểu thuộc dạng “Liên hệ tốc độ dài với tốc độ góc” ở tình huống số 4.
\choiceTF
{\True $v=\omega r$.}
{\True Cùng tốc độ góc, điểm xa tâm có tốc độ dài lớn hơn.}
{Tốc độ dài có đơn vị rad/s.}
{\True Các điểm trên vật rắn quay có cùng tốc độ góc.}
\loigiai{
\begin{itemchoice}
\item (Đúng) $v=\omega r$. (Vì): Phù hợp với định nghĩa hoặc hệ thức vật lí. b) Đúng: Phù hợp với định nghĩa hoặc hệ thức vật lí. c) Sai: Không phù hợp với định nghĩa hoặc hệ thức vật lí. d) Đúng: Phù hợp với định nghĩa hoặc hệ thức vật lí
\item (Đúng) Cùng tốc độ góc, điểm xa tâm có tốc độ dài lớn hơn.
\item (Sai) Tốc độ dài có đơn vị rad/s.
\item (Đúng) Các điểm trên vật rắn quay có cùng tốc độ góc.
\end{itemchoice}
}
\end{ex}

% ===== Câu 71 =====
% ID: AIQ_L10C6_FULL_0069
% Mức: VDC
\begin{ex}
% ID: AIQ_L10C6_FULL_0069
% Mức: VDC
Xét các phát biểu thuộc dạng “Liên hệ tốc độ dài với tốc độ góc” ở tình huống số 5.
\choiceTF
{\True $v=\omega r$.}
{\True Cùng tốc độ góc, điểm xa tâm có tốc độ dài lớn hơn.}
{Tốc độ dài có đơn vị rad/s.}
{\True Các điểm trên vật rắn quay có cùng tốc độ góc.}
\loigiai{
\begin{itemchoice}
\item (Đúng) $v=\omega r$. (Vì): Phù hợp với định nghĩa hoặc hệ thức vật lí. b) Đúng: Phù hợp với định nghĩa hoặc hệ thức vật lí. c) Sai: Không phù hợp với định nghĩa hoặc hệ thức vật lí. d) Đúng: Phù hợp với định nghĩa hoặc hệ thức vật lí
\item (Đúng) Cùng tốc độ góc, điểm xa tâm có tốc độ dài lớn hơn.
\item (Sai) Tốc độ dài có đơn vị rad/s.
\item (Đúng) Các điểm trên vật rắn quay có cùng tốc độ góc.
\end{itemchoice}
}
\end{ex}

% ===== Câu 72 =====
% ID: AIQ_L10C6_FULL_0070
% Mức: TH
\begin{ex}
% ID: AIQ_L10C6_FULL_0070
% Mức: TH
Một điểm quay với tốc độ góc $6$ rad/s trên quỹ đạo bán kính $0,2$ m. Tính tốc độ dài.
\shortans{1,2}
\loigiai{
$v=\omega r=6\cdot0,2$. Suy ra kết quả $1,2\ \mathrm{m/s}$. Đáp án trả lời ngắn không ghi đơn vị.
}
\end{ex}

% ===== Câu 73 =====
% ID: AIQ_L10C6_FULL_0071
% Mức: TH
\begin{ex}
% ID: AIQ_L10C6_FULL_0071
% Mức: TH
Một điểm quay với tốc độ góc $7$ rad/s trên quỹ đạo bán kính $0,3$ m. Tính tốc độ dài.
\shortans{2,1}
\loigiai{
$v=\omega r=7\cdot0,3$. Suy ra kết quả $2,1\ \mathrm{m/s}$. Đáp án trả lời ngắn không ghi đơn vị.
}
\end{ex}

% ===== Câu 74 =====
% ID: AIQ_L10C6_FULL_0072
% Mức: VD
\begin{ex}
% ID: AIQ_L10C6_FULL_0072
% Mức: VD
Một điểm quay với tốc độ góc $2$ rad/s trên quỹ đạo bán kính $0,4$ m. Tính tốc độ dài.
\shortans{0,8}
\loigiai{
$v=\omega r=2\cdot0,4$. Suy ra kết quả $0,8\ \mathrm{m/s}$. Đáp án trả lời ngắn không ghi đơn vị.
}
\end{ex}

% ===== Câu 75 =====
% ID: AIQ_L10C6_FULL_0073
% Mức: VD
\begin{ex}
% ID: AIQ_L10C6_FULL_0073
% Mức: VD
Một điểm quay với tốc độ góc $3$ rad/s trên quỹ đạo bán kính $0,5$ m. Tính tốc độ dài.
\shortans{1,5}
\loigiai{
$v=\omega r=3\cdot0,5$. Suy ra kết quả $1,5\ \mathrm{m/s}$. Đáp án trả lời ngắn không ghi đơn vị.
}
\end{ex}

% ===== Câu 76 =====
% ID: AIQ_L10C6_FULL_0074
% Mức: VDC
\begin{ex}
% ID: AIQ_L10C6_FULL_0074
% Mức: VDC
Một điểm quay với tốc độ góc $4$ rad/s trên quỹ đạo bán kính $0,6$ m. Tính tốc độ dài.
\shortans{2,4}
\loigiai{
$v=\omega r=4\cdot0,6$. Suy ra kết quả $2,4\ \mathrm{m/s}$. Đáp án trả lời ngắn không ghi đơn vị.
}
\end{ex}

% ===== Câu 77 =====
% ID: AIQ_L10C6_FULL_0075
% Mức: VDC
\begin{ex}
% ID: AIQ_L10C6_FULL_0075
% Mức: VDC
Một điểm quay với tốc độ góc $5$ rad/s trên quỹ đạo bán kính $0,2$ m. Tính tốc độ dài.
\shortans{1}
\loigiai{
$v=\omega r=5\cdot0,2$. Suy ra kết quả $1\ \mathrm{m/s}$. Đáp án trả lời ngắn không ghi đơn vị.
}
\end{ex}

% ===== Câu 78 =====
% ID: AIQ_L10C6_FULL_0076
% Mức: TH
\begin{bt}
% ID: AIQ_L10C6_FULL_0076
% Mức: TH
Trình bày cách giải: Một điểm quay với tốc độ góc $6$ rad/s trên quỹ đạo bán kính $0,3$ m. Tính tốc độ dài.
\loigiai{
$v=\omega r=6\cdot0,3$. Suy ra kết quả $1,8\ \mathrm{m/s}$.
}
\end{bt}

% ===== Câu 79 =====
% ID: AIQ_L10C6_FULL_0077
% Mức: TH
\begin{bt}
% ID: AIQ_L10C6_FULL_0077
% Mức: TH
Trình bày cách giải: Một điểm quay với tốc độ góc $7$ rad/s trên quỹ đạo bán kính $0,4$ m. Tính tốc độ dài.
\loigiai{
$v=\omega r=7\cdot0,4$. Suy ra kết quả $2,8\ \mathrm{m/s}$.
}
\end{bt}

% ===== Câu 80 =====
% ID: AIQ_L10C6_FULL_0078
% Mức: VD
\begin{bt}
% ID: AIQ_L10C6_FULL_0078
% Mức: VD
Trình bày cách giải: Một điểm quay với tốc độ góc $2$ rad/s trên quỹ đạo bán kính $0,5$ m. Tính tốc độ dài.
\loigiai{
$v=\omega r=2\cdot0,5$. Suy ra kết quả $1\ \mathrm{m/s}$.
}
\end{bt}

% ===== Câu 81 =====
% ID: AIQ_L10C6_FULL_0079
% Mức: VD
\begin{bt}
% ID: AIQ_L10C6_FULL_0079
% Mức: VD
Trình bày cách giải: Một điểm quay với tốc độ góc $3$ rad/s trên quỹ đạo bán kính $0,6$ m. Tính tốc độ dài.
\loigiai{
$v=\omega r=3\cdot0,6$. Suy ra kết quả $1,8\ \mathrm{m/s}$.
}
\end{bt}

% ===== Câu 82 =====
% ID: AIQ_L10C6_FULL_0080
% Mức: VDC
\begin{bt}
% ID: AIQ_L10C6_FULL_0080
% Mức: VDC
Trình bày cách giải: Một điểm quay với tốc độ góc $4$ rad/s trên quỹ đạo bán kính $0,2$ m. Tính tốc độ dài.
\loigiai{
$v=\omega r=4\cdot0,2$. Suy ra kết quả $0,8\ \mathrm{m/s}$.
}
\end{bt}

% ===== Câu 83 =====
% ID: AIQ_L10C6_FULL_0081
% Mức: VDC
\begin{bt}
% ID: AIQ_L10C6_FULL_0081
% Mức: VDC
Trình bày cách giải: Một điểm quay với tốc độ góc $5$ rad/s trên quỹ đạo bán kính $0,3$ m. Tính tốc độ dài.
\loigiai{
$v=\omega r=5\cdot0,3$. Suy ra kết quả $1,5\ \mathrm{m/s}$.
}
\end{bt}

% ===== Câu 84 =====
% ID: AIQ_L10C6_FULL_0082
% Mức: NB
\dangbt{Xác định góc quét, số vòng và thời gian chuyển động}
\begin{ex}
% ID: AIQ_L10C6_FULL_0082
% Mức: NB
Vật chuyển động tròn đều với tốc độ góc $1$ rad/s trong $2$ s. Tính góc quét.
\choice
{\True $2\ \mathrm{rad}$}
{$3\ \mathrm{rad}$}
{$0,5\ \mathrm{rad}$}
{Không đủ dữ kiện để có giá trị này}
\loigiai{
$\theta=\omega t=1\cdot2$. Suy ra kết quả $2\ \mathrm{rad}$.
}
\end{ex}

% ===== Câu 85 =====
% ID: AIQ_L10C6_FULL_0083
% Mức: NB
\begin{ex}
% ID: AIQ_L10C6_FULL_0083
% Mức: NB
Vật chuyển động tròn đều với tốc độ góc $2$ rad/s trong $3$ s. Tính góc quét.
\choice
{$5\ \mathrm{rad}$}
{\True $6\ \mathrm{rad}$}
{$0,67\ \mathrm{rad}$}
{$1,5\ \mathrm{rad}$}
\loigiai{
$\theta=\omega t=2\cdot3$. Suy ra kết quả $6\ \mathrm{rad}$.
}
\end{ex}

% ===== Câu 86 =====
% ID: AIQ_L10C6_FULL_0084
% Mức: TH
\begin{ex}
% ID: AIQ_L10C6_FULL_0084
% Mức: TH
Vật chuyển động tròn đều với tốc độ góc $3$ rad/s trong $4$ s. Tính góc quét.
\choice
{$7\ \mathrm{rad}$}
{$0,75\ \mathrm{rad}$}
{\True $12\ \mathrm{rad}$}
{$1,33\ \mathrm{rad}$}
\loigiai{
$\theta=\omega t=3\cdot4$. Suy ra kết quả $12\ \mathrm{rad}$.
}
\end{ex}

% ===== Câu 87 =====
% ID: AIQ_L10C6_FULL_0085
% Mức: TH
\begin{ex}
% ID: AIQ_L10C6_FULL_0085
% Mức: TH
Vật chuyển động tròn đều với tốc độ góc $4$ rad/s trong $5$ s. Tính góc quét.
\choice
{$9\ \mathrm{rad}$}
{$0,8\ \mathrm{rad}$}
{$1,25\ \mathrm{rad}$}
{\True $20\ \mathrm{rad}$}
\loigiai{
$\theta=\omega t=4\cdot5$. Suy ra kết quả $20\ \mathrm{rad}$.
}
\end{ex}

% ===== Câu 88 =====
% ID: AIQ_L10C6_FULL_0086
% Mức: VD
\begin{ex}
% ID: AIQ_L10C6_FULL_0086
% Mức: VD
Vật chuyển động tròn đều với tốc độ góc $5$ rad/s trong $6$ s. Tính góc quét.
\choice
{\True $30\ \mathrm{rad}$}
{$11\ \mathrm{rad}$}
{$0,83\ \mathrm{rad}$}
{$1,2\ \mathrm{rad}$}
\loigiai{
$\theta=\omega t=5\cdot6$. Suy ra kết quả $30\ \mathrm{rad}$.
}
\end{ex}

% ===== Câu 89 =====
% ID: AIQ_L10C6_FULL_0087
% Mức: VD
\begin{ex}
% ID: AIQ_L10C6_FULL_0087
% Mức: VD
Vật chuyển động tròn đều với tốc độ góc $1$ rad/s trong $7$ s. Tính góc quét.
\choice
{$8\ \mathrm{rad}$}
{\True $7\ \mathrm{rad}$}
{$0,14\ \mathrm{rad}$}
{Không đủ dữ kiện để có giá trị này}
\loigiai{
$\theta=\omega t=1\cdot7$. Suy ra kết quả $7\ \mathrm{rad}$.
}
\end{ex}

% ===== Câu 90 =====
% ID: AIQ_L10C6_FULL_0088
% Mức: VD
\begin{ex}
% ID: AIQ_L10C6_FULL_0088
% Mức: VD
Vật chuyển động tròn đều với tốc độ góc $2$ rad/s trong $8$ s. Tính góc quét.
\choice
{$10\ \mathrm{rad}$}
{$0,25\ \mathrm{rad}$}
{\True $16\ \mathrm{rad}$}
{$4\ \mathrm{rad}$}
\loigiai{
$\theta=\omega t=2\cdot8$. Suy ra kết quả $16\ \mathrm{rad}$.
}
\end{ex}

% ===== Câu 91 =====
% ID: AIQ_L10C6_FULL_0089
% Mức: VDC
\begin{ex}
% ID: AIQ_L10C6_FULL_0089
% Mức: VDC
Vật chuyển động tròn đều với tốc độ góc $3$ rad/s trong $9$ s. Tính góc quét.
\choice
{$12\ \mathrm{rad}$}
{$0,33\ \mathrm{rad}$}
{$3\ \mathrm{rad}$}
{\True $27\ \mathrm{rad}$}
\loigiai{
$\theta=\omega t=3\cdot9$. Suy ra kết quả $27\ \mathrm{rad}$.
}
\end{ex}

% ===== Câu 92 =====
% ID: AIQ_L10C6_FULL_0090
% Mức: VDC
\begin{ex}
% ID: AIQ_L10C6_FULL_0090
% Mức: VDC
Vật chuyển động tròn đều với tốc độ góc $4$ rad/s trong $10$ s. Tính góc quét.
\choice
{\True $40\ \mathrm{rad}$}
{$14\ \mathrm{rad}$}
{$0,4\ \mathrm{rad}$}
{$2,5\ \mathrm{rad}$}
\loigiai{
$\theta=\omega t=4\cdot10$. Suy ra kết quả $40\ \mathrm{rad}$.
}
\end{ex}

% ===== Câu 93 =====
% ID: AIQ_L10C6_FULL_0091
% Mức: VDC
\begin{ex}
% ID: AIQ_L10C6_FULL_0091
% Mức: VDC
Vật chuyển động tròn đều với tốc độ góc $5$ rad/s trong $11$ s. Tính góc quét.
\choice
{$16\ \mathrm{rad}$}
{\True $55\ \mathrm{rad}$}
{$0,45\ \mathrm{rad}$}
{$2,2\ \mathrm{rad}$}
\loigiai{
$\theta=\omega t=5\cdot11$. Suy ra kết quả $55\ \mathrm{rad}$.
}
\end{ex}

% ===== Câu 94 =====
% ID: AIQ_L10C6_FULL_0092
% Mức: NB
\begin{ex}
% ID: AIQ_L10C6_FULL_0092
% Mức: NB
Xét các phát biểu thuộc dạng “Xác định góc quét, số vòng và thời gian chuyển động” ở tình huống số 1.
\choiceTF
{\True Góc quét $\theta=\omega t$.}
{\True Số vòng $n=\theta/(2\pi)$.}
{Một vòng ứng với $\pi$ rad.}
{\True Chuyển động đều có góc quét tỉ lệ thời gian.}
\loigiai{
\begin{itemchoice}
\item (Đúng) Góc quét $\theta=\omega t$. (Vì): Phù hợp với định nghĩa hoặc hệ thức vật lí. b) Đúng: Phù hợp với định nghĩa hoặc hệ thức vật lí. c) Sai: Không phù hợp với định nghĩa hoặc hệ thức vật lí. d) Đúng: Phù hợp với định nghĩa hoặc hệ thức vật lí
\item (Đúng) Số vòng $n=\theta/(2\pi)$.
\item (Sai) Một vòng ứng với $\pi$ rad.
\item (Đúng) Chuyển động đều có góc quét tỉ lệ thời gian.
\end{itemchoice}
}
\end{ex}

% ===== Câu 95 =====
% ID: AIQ_L10C6_FULL_0093
% Mức: TH
\begin{ex}
% ID: AIQ_L10C6_FULL_0093
% Mức: TH
Xét các phát biểu thuộc dạng “Xác định góc quét, số vòng và thời gian chuyển động” ở tình huống số 2.
\choiceTF
{\True Góc quét $\theta=\omega t$.}
{\True Số vòng $n=\theta/(2\pi)$.}
{Một vòng ứng với $\pi$ rad.}
{\True Chuyển động đều có góc quét tỉ lệ thời gian.}
\loigiai{
\begin{itemchoice}
\item (Đúng) Góc quét $\theta=\omega t$. (Vì): Phù hợp với định nghĩa hoặc hệ thức vật lí. b) Đúng: Phù hợp với định nghĩa hoặc hệ thức vật lí. c) Sai: Không phù hợp với định nghĩa hoặc hệ thức vật lí. d) Đúng: Phù hợp với định nghĩa hoặc hệ thức vật lí
\item (Đúng) Số vòng $n=\theta/(2\pi)$.
\item (Sai) Một vòng ứng với $\pi$ rad.
\item (Đúng) Chuyển động đều có góc quét tỉ lệ thời gian.
\end{itemchoice}
}
\end{ex}

% ===== Câu 96 =====
% ID: AIQ_L10C6_FULL_0094
% Mức: TH
\begin{ex}
% ID: AIQ_L10C6_FULL_0094
% Mức: TH
Xét các phát biểu thuộc dạng “Xác định góc quét, số vòng và thời gian chuyển động” ở tình huống số 3.
\choiceTF
{\True Góc quét $\theta=\omega t$.}
{\True Số vòng $n=\theta/(2\pi)$.}
{Một vòng ứng với $\pi$ rad.}
{\True Chuyển động đều có góc quét tỉ lệ thời gian.}
\loigiai{
\begin{itemchoice}
\item (Đúng) Góc quét $\theta=\omega t$. (Vì): Phù hợp với định nghĩa hoặc hệ thức vật lí. b) Đúng: Phù hợp với định nghĩa hoặc hệ thức vật lí. c) Sai: Không phù hợp với định nghĩa hoặc hệ thức vật lí. d) Đúng: Phù hợp với định nghĩa hoặc hệ thức vật lí
\item (Đúng) Số vòng $n=\theta/(2\pi)$.
\item (Sai) Một vòng ứng với $\pi$ rad.
\item (Đúng) Chuyển động đều có góc quét tỉ lệ thời gian.
\end{itemchoice}
}
\end{ex}

% ===== Câu 97 =====
% ID: AIQ_L10C6_FULL_0095
% Mức: VD
\begin{ex}
% ID: AIQ_L10C6_FULL_0095
% Mức: VD
Xét các phát biểu thuộc dạng “Xác định góc quét, số vòng và thời gian chuyển động” ở tình huống số 4.
\choiceTF
{\True Góc quét $\theta=\omega t$.}
{\True Số vòng $n=\theta/(2\pi)$.}
{Một vòng ứng với $\pi$ rad.}
{\True Chuyển động đều có góc quét tỉ lệ thời gian.}
\loigiai{
\begin{itemchoice}
\item (Đúng) Góc quét $\theta=\omega t$. (Vì): Phù hợp với định nghĩa hoặc hệ thức vật lí. b) Đúng: Phù hợp với định nghĩa hoặc hệ thức vật lí. c) Sai: Không phù hợp với định nghĩa hoặc hệ thức vật lí. d) Đúng: Phù hợp với định nghĩa hoặc hệ thức vật lí
\item (Đúng) Số vòng $n=\theta/(2\pi)$.
\item (Sai) Một vòng ứng với $\pi$ rad.
\item (Đúng) Chuyển động đều có góc quét tỉ lệ thời gian.
\end{itemchoice}
}
\end{ex}

% ===== Câu 98 =====
% ID: AIQ_L10C6_FULL_0096
% Mức: VDC
\begin{ex}
% ID: AIQ_L10C6_FULL_0096
% Mức: VDC
Xét các phát biểu thuộc dạng “Xác định góc quét, số vòng và thời gian chuyển động” ở tình huống số 5.
\choiceTF
{\True Góc quét $\theta=\omega t$.}
{\True Số vòng $n=\theta/(2\pi)$.}
{Một vòng ứng với $\pi$ rad.}
{\True Chuyển động đều có góc quét tỉ lệ thời gian.}
\loigiai{
\begin{itemchoice}
\item (Đúng) Góc quét $\theta=\omega t$. (Vì): Phù hợp với định nghĩa hoặc hệ thức vật lí. b) Đúng: Phù hợp với định nghĩa hoặc hệ thức vật lí. c) Sai: Không phù hợp với định nghĩa hoặc hệ thức vật lí. d) Đúng: Phù hợp với định nghĩa hoặc hệ thức vật lí
\item (Đúng) Số vòng $n=\theta/(2\pi)$.
\item (Sai) Một vòng ứng với $\pi$ rad.
\item (Đúng) Chuyển động đều có góc quét tỉ lệ thời gian.
\end{itemchoice}
}
\end{ex}

% ===== Câu 99 =====
% ID: AIQ_L10C6_FULL_0097
% Mức: TH
\begin{ex}
% ID: AIQ_L10C6_FULL_0097
% Mức: TH
Vật chuyển động tròn đều với tốc độ góc $1$ rad/s trong $12$ s. Tính góc quét.
\shortans{12}
\loigiai{
$\theta=\omega t=1\cdot12$. Suy ra kết quả $12\ \mathrm{rad}$. Đáp án trả lời ngắn không ghi đơn vị.
}
\end{ex}

% ===== Câu 100 =====
% ID: AIQ_L10C6_FULL_0098
% Mức: TH
\begin{ex}
% ID: AIQ_L10C6_FULL_0098
% Mức: TH
Vật chuyển động tròn đều với tốc độ góc $2$ rad/s trong $13$ s. Tính góc quét.
\shortans{26}
\loigiai{
$\theta=\omega t=2\cdot13$. Suy ra kết quả $26\ \mathrm{rad}$. Đáp án trả lời ngắn không ghi đơn vị.
}
\end{ex}

% ===== Câu 101 =====
% ID: AIQ_L10C6_FULL_0099
% Mức: VD
\begin{ex}
% ID: AIQ_L10C6_FULL_0099
% Mức: VD
Vật chuyển động tròn đều với tốc độ góc $3$ rad/s trong $14$ s. Tính góc quét.
\shortans{42}
\loigiai{
$\theta=\omega t=3\cdot14$. Suy ra kết quả $42\ \mathrm{rad}$. Đáp án trả lời ngắn không ghi đơn vị.
}
\end{ex}

% ===== Câu 102 =====
% ID: AIQ_L10C6_FULL_0100
% Mức: VD
\begin{ex}
% ID: AIQ_L10C6_FULL_0100
% Mức: VD
Vật chuyển động tròn đều với tốc độ góc $4$ rad/s trong $15$ s. Tính góc quét.
\shortans{60}
\loigiai{
$\theta=\omega t=4\cdot15$. Suy ra kết quả $60\ \mathrm{rad}$. Đáp án trả lời ngắn không ghi đơn vị.
}
\end{ex}

% ===== Câu 103 =====
% ID: AIQ_L10C6_FULL_0101
% Mức: VDC
\begin{ex}
% ID: AIQ_L10C6_FULL_0101
% Mức: VDC
Vật chuyển động tròn đều với tốc độ góc $5$ rad/s trong $16$ s. Tính góc quét.
\shortans{80}
\loigiai{
$\theta=\omega t=5\cdot16$. Suy ra kết quả $80\ \mathrm{rad}$. Đáp án trả lời ngắn không ghi đơn vị.
}
\end{ex}

% ===== Câu 104 =====
% ID: AIQ_L10C6_FULL_0102
% Mức: VDC
\begin{ex}
% ID: AIQ_L10C6_FULL_0102
% Mức: VDC
Vật chuyển động tròn đều với tốc độ góc $1$ rad/s trong $17$ s. Tính góc quét.
\shortans{17}
\loigiai{
$\theta=\omega t=1\cdot17$. Suy ra kết quả $17\ \mathrm{rad}$. Đáp án trả lời ngắn không ghi đơn vị.
}
\end{ex}

% ===== Câu 105 =====
% ID: AIQ_L10C6_FULL_0103
% Mức: TH
\begin{bt}
% ID: AIQ_L10C6_FULL_0103
% Mức: TH
Trình bày cách giải: Vật chuyển động tròn đều với tốc độ góc $2$ rad/s trong $18$ s. Tính góc quét.
\loigiai{
$\theta=\omega t=2\cdot18$. Suy ra kết quả $36\ \mathrm{rad}$.
}
\end{bt}

% ===== Câu 106 =====
% ID: AIQ_L10C6_FULL_0104
% Mức: TH
\begin{bt}
% ID: AIQ_L10C6_FULL_0104
% Mức: TH
Trình bày cách giải: Vật chuyển động tròn đều với tốc độ góc $3$ rad/s trong $19$ s. Tính góc quét.
\loigiai{
$\theta=\omega t=3\cdot19$. Suy ra kết quả $57\ \mathrm{rad}$.
}
\end{bt}

% ===== Câu 107 =====
% ID: AIQ_L10C6_FULL_0105
% Mức: VD
\begin{bt}
% ID: AIQ_L10C6_FULL_0105
% Mức: VD
Trình bày cách giải: Vật chuyển động tròn đều với tốc độ góc $4$ rad/s trong $20$ s. Tính góc quét.
\loigiai{
$\theta=\omega t=4\cdot20$. Suy ra kết quả $80\ \mathrm{rad}$.
}
\end{bt}

% ===== Câu 108 =====
% ID: AIQ_L10C6_FULL_0106
% Mức: VD
\begin{bt}
% ID: AIQ_L10C6_FULL_0106
% Mức: VD
Trình bày cách giải: Vật chuyển động tròn đều với tốc độ góc $5$ rad/s trong $21$ s. Tính góc quét.
\loigiai{
$\theta=\omega t=5\cdot21$. Suy ra kết quả $105\ \mathrm{rad}$.
}
\end{bt}

% ===== Câu 109 =====
% ID: AIQ_L10C6_FULL_0107
% Mức: VDC
\begin{bt}
% ID: AIQ_L10C6_FULL_0107
% Mức: VDC
Trình bày cách giải: Vật chuyển động tròn đều với tốc độ góc $1$ rad/s trong $22$ s. Tính góc quét.
\loigiai{
$\theta=\omega t=1\cdot22$. Suy ra kết quả $22\ \mathrm{rad}$.
}
\end{bt}

% ===== Câu 110 =====
% ID: AIQ_L10C6_FULL_0108
% Mức: VDC
\begin{bt}
% ID: AIQ_L10C6_FULL_0108
% Mức: VDC
Trình bày cách giải: Vật chuyển động tròn đều với tốc độ góc $2$ rad/s trong $23$ s. Tính góc quét.
\loigiai{
$\theta=\omega t=2\cdot23$. Suy ra kết quả $46\ \mathrm{rad}$.
}
\end{bt}

% ===== Câu 111 =====
% ID: AIQ_L10C6_FULL_0109
% Mức: NB
\dangbt{Bài toán truyền chuyển động tròn}
\begin{ex}
% ID: AIQ_L10C6_FULL_0109
% Mức: NB
Hai bánh xe tiếp xúc không trượt có bán kính $r_1=0,1$ m, $r_2=0,2$ m. Bánh 1 quay với $\omega_1=4$ rad/s. Tính độ lớn $\omega_2$.
\choice
{\True $2\ \mathrm{rad/s}$}
{$8\ \mathrm{rad/s}$}
{$0,4\ \mathrm{rad/s}$}
{$4\ \mathrm{rad/s}$}
\loigiai{
Không trượt: $\omega_1r_1=\omega_2r_2\Rightarrow\omega_2=4\cdot0,1/0,2$. Suy ra kết quả $2\ \mathrm{rad/s}$.
}
\end{ex}

% ===== Câu 112 =====
% ID: AIQ_L10C6_FULL_0110
% Mức: NB
\begin{ex}
% ID: AIQ_L10C6_FULL_0110
% Mức: NB
Hai bánh xe tiếp xúc không trượt có bán kính $r_1=0,12$ m, $r_2=0,23$ m. Bánh 1 quay với $\omega_1=5$ rad/s. Tính độ lớn $\omega_2$.
\choice
{$9,58\ \mathrm{rad/s}$}
{\True $2,61\ \mathrm{rad/s}$}
{$0,6\ \mathrm{rad/s}$}
{$5,22\ \mathrm{rad/s}$}
\loigiai{
Không trượt: $\omega_1r_1=\omega_2r_2\Rightarrow\omega_2=5\cdot0,12/0,23$. Suy ra kết quả $2,61\ \mathrm{rad/s}$.
}
\end{ex}

% ===== Câu 113 =====
% ID: AIQ_L10C6_FULL_0111
% Mức: TH
\begin{ex}
% ID: AIQ_L10C6_FULL_0111
% Mức: TH
Hai bánh xe tiếp xúc không trượt có bán kính $r_1=0,14$ m, $r_2=0,26$ m. Bánh 1 quay với $\omega_1=6$ rad/s. Tính độ lớn $\omega_2$.
\choice
{$11,14\ \mathrm{rad/s}$}
{$0,84\ \mathrm{rad/s}$}
{\True $3,23\ \mathrm{rad/s}$}
{$6,46\ \mathrm{rad/s}$}
\loigiai{
Không trượt: $\omega_1r_1=\omega_2r_2\Rightarrow\omega_2=6\cdot0,14/0,26$. Suy ra kết quả $3,23\ \mathrm{rad/s}$.
}
\end{ex}

% ===== Câu 114 =====
% ID: AIQ_L10C6_FULL_0112
% Mức: TH
\begin{ex}
% ID: AIQ_L10C6_FULL_0112
% Mức: TH
Hai bánh xe tiếp xúc không trượt có bán kính $r_1=0,16$ m, $r_2=0,29$ m. Bánh 1 quay với $\omega_1=7$ rad/s. Tính độ lớn $\omega_2$.
\choice
{$12,69\ \mathrm{rad/s}$}
{$1,12\ \mathrm{rad/s}$}
{$7,72\ \mathrm{rad/s}$}
{\True $3,86\ \mathrm{rad/s}$}
\loigiai{
Không trượt: $\omega_1r_1=\omega_2r_2\Rightarrow\omega_2=7\cdot0,16/0,29$. Suy ra kết quả $3,86\ \mathrm{rad/s}$.
}
\end{ex}

% ===== Câu 115 =====
% ID: AIQ_L10C6_FULL_0113
% Mức: VD
\begin{ex}
% ID: AIQ_L10C6_FULL_0113
% Mức: VD
Hai bánh xe tiếp xúc không trượt có bán kính $r_1=0,18$ m, $r_2=0,32$ m. Bánh 1 quay với $\omega_1=8$ rad/s. Tính độ lớn $\omega_2$.
\choice
{\True $4,5\ \mathrm{rad/s}$}
{$14,22\ \mathrm{rad/s}$}
{$1,44\ \mathrm{rad/s}$}
{$9\ \mathrm{rad/s}$}
\loigiai{
Không trượt: $\omega_1r_1=\omega_2r_2\Rightarrow\omega_2=8\cdot0,18/0,32$. Suy ra kết quả $4,5\ \mathrm{rad/s}$.
}
\end{ex}

% ===== Câu 116 =====
% ID: AIQ_L10C6_FULL_0114
% Mức: VD
\begin{ex}
% ID: AIQ_L10C6_FULL_0114
% Mức: VD
Hai bánh xe tiếp xúc không trượt có bán kính $r_1=0,2$ m, $r_2=0,35$ m. Bánh 1 quay với $\omega_1=4$ rad/s. Tính độ lớn $\omega_2$.
\choice
{$7\ \mathrm{rad/s}$}
{\True $2,29\ \mathrm{rad/s}$}
{$0,8\ \mathrm{rad/s}$}
{$4,57\ \mathrm{rad/s}$}
\loigiai{
Không trượt: $\omega_1r_1=\omega_2r_2\Rightarrow\omega_2=4\cdot0,2/0,35$. Suy ra kết quả $2,29\ \mathrm{rad/s}$.
}
\end{ex}

% ===== Câu 117 =====
% ID: AIQ_L10C6_FULL_0115
% Mức: VD
\begin{ex}
% ID: AIQ_L10C6_FULL_0115
% Mức: VD
Hai bánh xe tiếp xúc không trượt có bán kính $r_1=0,22$ m, $r_2=0,38$ m. Bánh 1 quay với $\omega_1=5$ rad/s. Tính độ lớn $\omega_2$.
\choice
{$8,64\ \mathrm{rad/s}$}
{$1,1\ \mathrm{rad/s}$}
{\True $2,89\ \mathrm{rad/s}$}
{$5,79\ \mathrm{rad/s}$}
\loigiai{
Không trượt: $\omega_1r_1=\omega_2r_2\Rightarrow\omega_2=5\cdot0,22/0,38$. Suy ra kết quả $2,89\ \mathrm{rad/s}$.
}
\end{ex}

% ===== Câu 118 =====
% ID: AIQ_L10C6_FULL_0116
% Mức: VDC
\begin{ex}
% ID: AIQ_L10C6_FULL_0116
% Mức: VDC
Hai bánh xe tiếp xúc không trượt có bán kính $r_1=0,24$ m, $r_2=0,41$ m. Bánh 1 quay với $\omega_1=6$ rad/s. Tính độ lớn $\omega_2$.
\choice
{$10,25\ \mathrm{rad/s}$}
{$1,44\ \mathrm{rad/s}$}
{$7,02\ \mathrm{rad/s}$}
{\True $3,51\ \mathrm{rad/s}$}
\loigiai{
Không trượt: $\omega_1r_1=\omega_2r_2\Rightarrow\omega_2=6\cdot0,24/0,41$. Suy ra kết quả $3,51\ \mathrm{rad/s}$.
}
\end{ex}

% ===== Câu 119 =====
% ID: AIQ_L10C6_FULL_0117
% Mức: VDC
\begin{ex}
% ID: AIQ_L10C6_FULL_0117
% Mức: VDC
Hai bánh xe tiếp xúc không trượt có bán kính $r_1=0,26$ m, $r_2=0,44$ m. Bánh 1 quay với $\omega_1=7$ rad/s. Tính độ lớn $\omega_2$.
\choice
{\True $4,14\ \mathrm{rad/s}$}
{$11,85\ \mathrm{rad/s}$}
{$1,82\ \mathrm{rad/s}$}
{$8,27\ \mathrm{rad/s}$}
\loigiai{
Không trượt: $\omega_1r_1=\omega_2r_2\Rightarrow\omega_2=7\cdot0,26/0,44$. Suy ra kết quả $4,14\ \mathrm{rad/s}$.
}
\end{ex}

% ===== Câu 120 =====
% ID: AIQ_L10C6_FULL_0118
% Mức: VDC
\begin{ex}
% ID: AIQ_L10C6_FULL_0118
% Mức: VDC
Hai bánh xe tiếp xúc không trượt có bán kính $r_1=0,28$ m, $r_2=0,47$ m. Bánh 1 quay với $\omega_1=8$ rad/s. Tính độ lớn $\omega_2$.
\choice
{$13,43\ \mathrm{rad/s}$}
{\True $4,77\ \mathrm{rad/s}$}
{$2,24\ \mathrm{rad/s}$}
{$9,53\ \mathrm{rad/s}$}
\loigiai{
Không trượt: $\omega_1r_1=\omega_2r_2\Rightarrow\omega_2=8\cdot0,28/0,47$. Suy ra kết quả $4,77\ \mathrm{rad/s}$.
}
\end{ex}

% ===== Câu 121 =====
% ID: AIQ_L10C6_FULL_0119
% Mức: NB
\begin{ex}
% ID: AIQ_L10C6_FULL_0119
% Mức: NB
Xét các phát biểu thuộc dạng “Bài toán truyền chuyển động tròn” ở tình huống số 1.
\choiceTF
{\True Hai bánh tiếp xúc không trượt có tốc độ dài tại tiếp điểm bằng nhau.}
{\True Bánh nhỏ thường quay nhanh hơn bánh lớn.}
{Hai bánh tiếp xúc quay cùng chiều.}
{\True Dùng $\omega_1r_1=\omega_2r_2$.}
\loigiai{
\begin{itemchoice}
\item (Đúng) Hai bánh tiếp xúc không trượt có tốc độ dài tại tiếp điểm bằng nhau. (Vì): Phù hợp với định nghĩa hoặc hệ thức vật lí. b) Đúng: Phù hợp với định nghĩa hoặc hệ thức vật lí. c) Sai: Không phù hợp với định nghĩa hoặc hệ thức vật lí. d) Đúng: Phù hợp với định nghĩa hoặc hệ thức vật lí
\item (Đúng) Bánh nhỏ thường quay nhanh hơn bánh lớn.
\item (Sai) Hai bánh tiếp xúc quay cùng chiều.
\item (Đúng) Dùng $\omega_1r_1=\omega_2r_2$.
\end{itemchoice}
}
\end{ex}

% ===== Câu 122 =====
% ID: AIQ_L10C6_FULL_0120
% Mức: TH
\begin{ex}
% ID: AIQ_L10C6_FULL_0120
% Mức: TH
Xét các phát biểu thuộc dạng “Bài toán truyền chuyển động tròn” ở tình huống số 2.
\choiceTF
{\True Hai bánh tiếp xúc không trượt có tốc độ dài tại tiếp điểm bằng nhau.}
{\True Bánh nhỏ thường quay nhanh hơn bánh lớn.}
{Hai bánh tiếp xúc quay cùng chiều.}
{\True Dùng $\omega_1r_1=\omega_2r_2$.}
\loigiai{
\begin{itemchoice}
\item (Đúng) Hai bánh tiếp xúc không trượt có tốc độ dài tại tiếp điểm bằng nhau. (Vì): Phù hợp với định nghĩa hoặc hệ thức vật lí. b) Đúng: Phù hợp với định nghĩa hoặc hệ thức vật lí. c) Sai: Không phù hợp với định nghĩa hoặc hệ thức vật lí. d) Đúng: Phù hợp với định nghĩa hoặc hệ thức vật lí
\item (Đúng) Bánh nhỏ thường quay nhanh hơn bánh lớn.
\item (Sai) Hai bánh tiếp xúc quay cùng chiều.
\item (Đúng) Dùng $\omega_1r_1=\omega_2r_2$.
\end{itemchoice}
}
\end{ex}

% ===== Câu 123 =====
% ID: AIQ_L10C6_FULL_0121
% Mức: TH
\begin{ex}
% ID: AIQ_L10C6_FULL_0121
% Mức: TH
Xét các phát biểu thuộc dạng “Bài toán truyền chuyển động tròn” ở tình huống số 3.
\choiceTF
{\True Hai bánh tiếp xúc không trượt có tốc độ dài tại tiếp điểm bằng nhau.}
{\True Bánh nhỏ thường quay nhanh hơn bánh lớn.}
{Hai bánh tiếp xúc quay cùng chiều.}
{\True Dùng $\omega_1r_1=\omega_2r_2$.}
\loigiai{
\begin{itemchoice}
\item (Đúng) Hai bánh tiếp xúc không trượt có tốc độ dài tại tiếp điểm bằng nhau. (Vì): Phù hợp với định nghĩa hoặc hệ thức vật lí. b) Đúng: Phù hợp với định nghĩa hoặc hệ thức vật lí. c) Sai: Không phù hợp với định nghĩa hoặc hệ thức vật lí. d) Đúng: Phù hợp với định nghĩa hoặc hệ thức vật lí
\item (Đúng) Bánh nhỏ thường quay nhanh hơn bánh lớn.
\item (Sai) Hai bánh tiếp xúc quay cùng chiều.
\item (Đúng) Dùng $\omega_1r_1=\omega_2r_2$.
\end{itemchoice}
}
\end{ex}

% ===== Câu 124 =====
% ID: AIQ_L10C6_FULL_0122
% Mức: VD
\begin{ex}
% ID: AIQ_L10C6_FULL_0122
% Mức: VD
Xét các phát biểu thuộc dạng “Bài toán truyền chuyển động tròn” ở tình huống số 4.
\choiceTF
{\True Hai bánh tiếp xúc không trượt có tốc độ dài tại tiếp điểm bằng nhau.}
{\True Bánh nhỏ thường quay nhanh hơn bánh lớn.}
{Hai bánh tiếp xúc quay cùng chiều.}
{\True Dùng $\omega_1r_1=\omega_2r_2$.}
\loigiai{
\begin{itemchoice}
\item (Đúng) Hai bánh tiếp xúc không trượt có tốc độ dài tại tiếp điểm bằng nhau. (Vì): Phù hợp với định nghĩa hoặc hệ thức vật lí. b) Đúng: Phù hợp với định nghĩa hoặc hệ thức vật lí. c) Sai: Không phù hợp với định nghĩa hoặc hệ thức vật lí. d) Đúng: Phù hợp với định nghĩa hoặc hệ thức vật lí
\item (Đúng) Bánh nhỏ thường quay nhanh hơn bánh lớn.
\item (Sai) Hai bánh tiếp xúc quay cùng chiều.
\item (Đúng) Dùng $\omega_1r_1=\omega_2r_2$.
\end{itemchoice}
}
\end{ex}

% ===== Câu 125 =====
% ID: AIQ_L10C6_FULL_0123
% Mức: VDC
\begin{ex}
% ID: AIQ_L10C6_FULL_0123
% Mức: VDC
Xét các phát biểu thuộc dạng “Bài toán truyền chuyển động tròn” ở tình huống số 5.
\choiceTF
{\True Hai bánh tiếp xúc không trượt có tốc độ dài tại tiếp điểm bằng nhau.}
{\True Bánh nhỏ thường quay nhanh hơn bánh lớn.}
{Hai bánh tiếp xúc quay cùng chiều.}
{\True Dùng $\omega_1r_1=\omega_2r_2$.}
\loigiai{
\begin{itemchoice}
\item (Đúng) Hai bánh tiếp xúc không trượt có tốc độ dài tại tiếp điểm bằng nhau. (Vì): Phù hợp với định nghĩa hoặc hệ thức vật lí. b) Đúng: Phù hợp với định nghĩa hoặc hệ thức vật lí. c) Sai: Không phù hợp với định nghĩa hoặc hệ thức vật lí. d) Đúng: Phù hợp với định nghĩa hoặc hệ thức vật lí
\item (Đúng) Bánh nhỏ thường quay nhanh hơn bánh lớn.
\item (Sai) Hai bánh tiếp xúc quay cùng chiều.
\item (Đúng) Dùng $\omega_1r_1=\omega_2r_2$.
\end{itemchoice}
}
\end{ex}

% ===== Câu 126 =====
% ID: AIQ_L10C6_FULL_0124
% Mức: TH
\begin{ex}
% ID: AIQ_L10C6_FULL_0124
% Mức: TH
Hai bánh xe tiếp xúc không trượt có bán kính $r_1=0,3$ m, $r_2=0,5$ m. Bánh 1 quay với $\omega_1=4$ rad/s. Tính độ lớn $\omega_2$.
\shortans{2,4}
\loigiai{
Không trượt: $\omega_1r_1=\omega_2r_2\Rightarrow\omega_2=4\cdot0,3/0,5$. Suy ra kết quả $2,4\ \mathrm{rad/s}$. Đáp án trả lời ngắn không ghi đơn vị.
}
\end{ex}

% ===== Câu 127 =====
% ID: AIQ_L10C6_FULL_0125
% Mức: TH
\begin{ex}
% ID: AIQ_L10C6_FULL_0125
% Mức: TH
Hai bánh xe tiếp xúc không trượt có bán kính $r_1=0,32$ m, $r_2=0,53$ m. Bánh 1 quay với $\omega_1=5$ rad/s. Tính độ lớn $\omega_2$.
\shortans{3,02}
\loigiai{
Không trượt: $\omega_1r_1=\omega_2r_2\Rightarrow\omega_2=5\cdot0,32/0,53$. Suy ra kết quả $3,02\ \mathrm{rad/s}$. Đáp án trả lời ngắn không ghi đơn vị.
}
\end{ex}

% ===== Câu 128 =====
% ID: AIQ_L10C6_FULL_0126
% Mức: VD
\begin{ex}
% ID: AIQ_L10C6_FULL_0126
% Mức: VD
Hai bánh xe tiếp xúc không trượt có bán kính $r_1=0,34$ m, $r_2=0,56$ m. Bánh 1 quay với $\omega_1=6$ rad/s. Tính độ lớn $\omega_2$.
\shortans{3,64}
\loigiai{
Không trượt: $\omega_1r_1=\omega_2r_2\Rightarrow\omega_2=6\cdot0,34/0,56$. Suy ra kết quả $3,64\ \mathrm{rad/s}$. Đáp án trả lời ngắn không ghi đơn vị.
}
\end{ex}

% ===== Câu 129 =====
% ID: AIQ_L10C6_FULL_0127
% Mức: VD
\begin{ex}
% ID: AIQ_L10C6_FULL_0127
% Mức: VD
Hai bánh xe tiếp xúc không trượt có bán kính $r_1=0,36$ m, $r_2=0,59$ m. Bánh 1 quay với $\omega_1=7$ rad/s. Tính độ lớn $\omega_2$.
\shortans{4,27}
\loigiai{
Không trượt: $\omega_1r_1=\omega_2r_2\Rightarrow\omega_2=7\cdot0,36/0,59$. Suy ra kết quả $4,27\ \mathrm{rad/s}$. Đáp án trả lời ngắn không ghi đơn vị.
}
\end{ex}

% ===== Câu 130 =====
% ID: AIQ_L10C6_FULL_0128
% Mức: VDC
\begin{ex}
% ID: AIQ_L10C6_FULL_0128
% Mức: VDC
Hai bánh xe tiếp xúc không trượt có bán kính $r_1=0,38$ m, $r_2=0,62$ m. Bánh 1 quay với $\omega_1=8$ rad/s. Tính độ lớn $\omega_2$.
\shortans{4,9}
\loigiai{
Không trượt: $\omega_1r_1=\omega_2r_2\Rightarrow\omega_2=8\cdot0,38/0,62$. Suy ra kết quả $4,9\ \mathrm{rad/s}$. Đáp án trả lời ngắn không ghi đơn vị.
}
\end{ex}

% ===== Câu 131 =====
% ID: AIQ_L10C6_FULL_0129
% Mức: VDC
\begin{ex}
% ID: AIQ_L10C6_FULL_0129
% Mức: VDC
Hai bánh xe tiếp xúc không trượt có bán kính $r_1=0,4$ m, $r_2=0,65$ m. Bánh 1 quay với $\omega_1=4$ rad/s. Tính độ lớn $\omega_2$.
\shortans{2,46}
\loigiai{
Không trượt: $\omega_1r_1=\omega_2r_2\Rightarrow\omega_2=4\cdot0,4/0,65$. Suy ra kết quả $2,46\ \mathrm{rad/s}$. Đáp án trả lời ngắn không ghi đơn vị.
}
\end{ex}

% ===== Câu 132 =====
% ID: AIQ_L10C6_FULL_0130
% Mức: TH
\begin{ex}
% ID: AIQ_L10C6_FULL_0130
% Mức: TH
Trình bày cách giải: Hai bánh xe tiếp xúc không trượt có bán kính $r_1=0,42$ m, $r_2=0,68$ m. Bánh 1 quay với $\omega_1=5$ rad/s. Tính độ lớn $\omega_2$.
\choiceTF

\loigiai{
Không trượt: $\omega_1r_1=\omega_2r_2\Rightarrow\omega_2=5\cdot0,42/0,68$. Suy ra kết quả $3,09\ \mathrm{rad/s}$.
}
\end{ex}

% ===== Câu 133 =====
% ID: AIQ_L10C6_FULL_0131
% Mức: TH
\begin{ex}
% ID: AIQ_L10C6_FULL_0131
% Mức: TH
Trình bày cách giải: Hai bánh xe tiếp xúc không trượt có bán kính $r_1=0,44$ m, $r_2=0,71$ m. Bánh 1 quay với $\omega_1=6$ rad/s. Tính độ lớn $\omega_2$.
\choiceTF

\loigiai{
Không trượt: $\omega_1r_1=\omega_2r_2\Rightarrow\omega_2=6\cdot0,44/0,71$. Suy ra kết quả $3,72\ \mathrm{rad/s}$.
}
\end{ex}

% ===== Câu 134 =====
% ID: AIQ_L10C6_FULL_0132
% Mức: VD
\begin{ex}
% ID: AIQ_L10C6_FULL_0132
% Mức: VD
Trình bày cách giải: Hai bánh xe tiếp xúc không trượt có bán kính $r_1=0,46$ m, $r_2=0,74$ m. Bánh 1 quay với $\omega_1=7$ rad/s. Tính độ lớn $\omega_2$.
\choiceTF

\loigiai{
Không trượt: $\omega_1r_1=\omega_2r_2\Rightarrow\omega_2=7\cdot0,46/0,74$. Suy ra kết quả $4,35\ \mathrm{rad/s}$.
}
\end{ex}

% ===== Câu 135 =====
% ID: AIQ_L10C6_FULL_0133
% Mức: VD
\begin{ex}
% ID: AIQ_L10C6_FULL_0133
% Mức: VD
Trình bày cách giải: Hai bánh xe tiếp xúc không trượt có bán kính $r_1=0,48$ m, $r_2=0,77$ m. Bánh 1 quay với $\omega_1=8$ rad/s. Tính độ lớn $\omega_2$.
\choiceTF

\loigiai{
Không trượt: $\omega_1r_1=\omega_2r_2\Rightarrow\omega_2=8\cdot0,48/0,77$. Suy ra kết quả $4,99\ \mathrm{rad/s}$.
}
\end{ex}

% ===== Câu 136 =====
% ID: AIQ_L10C6_FULL_0134
% Mức: VDC
\begin{ex}
% ID: AIQ_L10C6_FULL_0134
% Mức: VDC
Trình bày cách giải: Hai bánh xe tiếp xúc không trượt có bán kính $r_1=0,5$ m, $r_2=0,8$ m. Bánh 1 quay với $\omega_1=4$ rad/s. Tính độ lớn $\omega_2$.
\choiceTF

\loigiai{
Không trượt: $\omega_1r_1=\omega_2r_2\Rightarrow\omega_2=4\cdot0,5/0,8$. Suy ra kết quả $2,5\ \mathrm{rad/s}$.
}
\end{ex}

% ===== Câu 137 =====
% ID: AIQ_L10C6_FULL_0135
% Mức: VDC
\begin{ex}
% ID: AIQ_L10C6_FULL_0135
% Mức: VDC
Trình bày cách giải: Hai bánh xe tiếp xúc không trượt có bán kính $r_1=0,52$ m, $r_2=0,83$ m. Bánh 1 quay với $\omega_1=5$ rad/s. Tính độ lớn $\omega_2$.
\choiceTF

\loigiai{
Không trượt: $\omega_1r_1=\omega_2r_2\Rightarrow\omega_2=5\cdot0,52/0,83$. Suy ra kết quả $3,13\ \mathrm{rad/s}$.
}
\end{ex}

% ===== Câu 138 =====
% ID: AIQ_L10C6_FULL_0136
% Mức: NB
\dangbt{Khai thác đồ thị và tình huống thực tiễn về chuyển động tròn đều}
\begin{ex}
% ID: AIQ_L10C6_FULL_0136
% Mức: NB
Đầu kim giây dài $0,1$ m quay đều. Tính tốc độ dài của đầu kim.
\choice
{\True $0,01\ \mathrm{m/s}$}
{$0,63\ \mathrm{m/s}$}
{$0\ \mathrm{m/s}$}
{$6\ \mathrm{m/s}$}
\loigiai{
Kim giây có $T=60$ s, nên $v=2\pi r/T=2\pi\cdot0,1/60$. Suy ra kết quả $0,01\ \mathrm{m/s}$.
}
\end{ex}

% ===== Câu 139 =====
% ID: AIQ_L10C6_FULL_0137
% Mức: NB
\begin{ex}
% ID: AIQ_L10C6_FULL_0137
% Mức: NB
Đầu kim giây dài $0,11$ m quay đều. Tính tốc độ dài của đầu kim.
\choice
{$0,69\ \mathrm{m/s}$}
{\True $0,01\ \mathrm{m/s}$}
{$0\ \mathrm{m/s}$}
{$6,6\ \mathrm{m/s}$}
\loigiai{
Kim giây có $T=60$ s, nên $v=2\pi r/T=2\pi\cdot0,11/60$. Suy ra kết quả $0,01\ \mathrm{m/s}$.
}
\end{ex}

% ===== Câu 140 =====
% ID: AIQ_L10C6_FULL_0138
% Mức: TH
\begin{ex}
% ID: AIQ_L10C6_FULL_0138
% Mức: TH
Đầu kim giây dài $0,12$ m quay đều. Tính tốc độ dài của đầu kim.
\choice
{$0,75\ \mathrm{m/s}$}
{$0\ \mathrm{m/s}$}
{\True $0,01\ \mathrm{m/s}$}
{$7,2\ \mathrm{m/s}$}
\loigiai{
Kim giây có $T=60$ s, nên $v=2\pi r/T=2\pi\cdot0,12/60$. Suy ra kết quả $0,01\ \mathrm{m/s}$.
}
\end{ex}

% ===== Câu 141 =====
% ID: AIQ_L10C6_FULL_0139
% Mức: TH
\begin{ex}
% ID: AIQ_L10C6_FULL_0139
% Mức: TH
Đầu kim giây dài $0,13$ m quay đều. Tính tốc độ dài của đầu kim.
\choice
{$0,82\ \mathrm{m/s}$}
{$0\ \mathrm{m/s}$}
{$7,8\ \mathrm{m/s}$}
{\True $0,01\ \mathrm{m/s}$}
\loigiai{
Kim giây có $T=60$ s, nên $v=2\pi r/T=2\pi\cdot0,13/60$. Suy ra kết quả $0,01\ \mathrm{m/s}$.
}
\end{ex}

% ===== Câu 142 =====
% ID: AIQ_L10C6_FULL_0140
% Mức: VD
\begin{ex}
% ID: AIQ_L10C6_FULL_0140
% Mức: VD
Đầu kim giây dài $0,14$ m quay đều. Tính tốc độ dài của đầu kim.
\choice
{\True $0,01\ \mathrm{m/s}$}
{$0,88\ \mathrm{m/s}$}
{$0\ \mathrm{m/s}$}
{$8,4\ \mathrm{m/s}$}
\loigiai{
Kim giây có $T=60$ s, nên $v=2\pi r/T=2\pi\cdot0,14/60$. Suy ra kết quả $0,01\ \mathrm{m/s}$.
}
\end{ex}

% ===== Câu 143 =====
% ID: AIQ_L10C6_FULL_0141
% Mức: VD
\begin{ex}
% ID: AIQ_L10C6_FULL_0141
% Mức: VD
Đầu kim giây dài $0,15$ m quay đều. Tính tốc độ dài của đầu kim.
\choice
{$0,94\ \mathrm{m/s}$}
{\True $0,02\ \mathrm{m/s}$}
{$0\ \mathrm{m/s}$}
{$9\ \mathrm{m/s}$}
\loigiai{
Kim giây có $T=60$ s, nên $v=2\pi r/T=2\pi\cdot0,15/60$. Suy ra kết quả $0,02\ \mathrm{m/s}$.
}
\end{ex}

% ===== Câu 144 =====
% ID: AIQ_L10C6_FULL_0142
% Mức: VD
\begin{ex}
% ID: AIQ_L10C6_FULL_0142
% Mức: VD
Đầu kim giây dài $0,16$ m quay đều. Tính tốc độ dài của đầu kim.
\choice
{$1,01\ \mathrm{m/s}$}
{$0\ \mathrm{m/s}$}
{\True $0,02\ \mathrm{m/s}$}
{$9,6\ \mathrm{m/s}$}
\loigiai{
Kim giây có $T=60$ s, nên $v=2\pi r/T=2\pi\cdot0,16/60$. Suy ra kết quả $0,02\ \mathrm{m/s}$.
}
\end{ex}

% ===== Câu 145 =====
% ID: AIQ_L10C6_FULL_0143
% Mức: VDC
\begin{ex}
% ID: AIQ_L10C6_FULL_0143
% Mức: VDC
Đầu kim giây dài $0,17$ m quay đều. Tính tốc độ dài của đầu kim.
\choice
{$1,07\ \mathrm{m/s}$}
{$0\ \mathrm{m/s}$}
{$10,2\ \mathrm{m/s}$}
{\True $0,02\ \mathrm{m/s}$}
\loigiai{
Kim giây có $T=60$ s, nên $v=2\pi r/T=2\pi\cdot0,17/60$. Suy ra kết quả $0,02\ \mathrm{m/s}$.
}
\end{ex}

% ===== Câu 146 =====
% ID: AIQ_L10C6_FULL_0144
% Mức: VDC
\begin{ex}
% ID: AIQ_L10C6_FULL_0144
% Mức: VDC
Đầu kim giây dài $0,18$ m quay đều. Tính tốc độ dài của đầu kim.
\choice
{\True $0,02\ \mathrm{m/s}$}
{$1,13\ \mathrm{m/s}$}
{$0\ \mathrm{m/s}$}
{$10,8\ \mathrm{m/s}$}
\loigiai{
Kim giây có $T=60$ s, nên $v=2\pi r/T=2\pi\cdot0,18/60$. Suy ra kết quả $0,02\ \mathrm{m/s}$.
}
\end{ex}

% ===== Câu 147 =====
% ID: AIQ_L10C6_FULL_0145
% Mức: VDC
\begin{ex}
% ID: AIQ_L10C6_FULL_0145
% Mức: VDC
Đầu kim giây dài $0,19$ m quay đều. Tính tốc độ dài của đầu kim.
\choice
{$1,19\ \mathrm{m/s}$}
{\True $0,02\ \mathrm{m/s}$}
{$0\ \mathrm{m/s}$}
{$11,4\ \mathrm{m/s}$}
\loigiai{
Kim giây có $T=60$ s, nên $v=2\pi r/T=2\pi\cdot0,19/60$. Suy ra kết quả $0,02\ \mathrm{m/s}$.
}
\end{ex}

% ===== Câu 148 =====
% ID: AIQ_L10C6_FULL_0146
% Mức: NB
\begin{ex}
% ID: AIQ_L10C6_FULL_0146
% Mức: NB
Xét các phát biểu thuộc dạng “Khai thác đồ thị và tình huống thực tiễn về chuyển động tròn đều” ở tình huống số 1.
\choiceTF
{\True Kim giây có chu kì $60\,\text{s}$.}
{\True Đầu kim dài hơn có tốc độ dài lớn hơn nếu cùng tốc độ góc.}
{Mọi kim đồng hồ có cùng chu kì.}
{\True Đồ thị góc theo thời gian của chuyển động tròn đều là đường thẳng.}
\loigiai{
\begin{itemchoice}
\item (Đúng) Kim giây có chu kì $60\,\text{s}$. (Vì): Phù hợp với định nghĩa hoặc hệ thức vật lí. b) Đúng: Phù hợp với định nghĩa hoặc hệ thức vật lí. c) Sai: Không phù hợp với định nghĩa hoặc hệ thức vật lí. d) Đúng: Phù hợp với định nghĩa hoặc hệ thức vật lí
\item (Đúng) Đầu kim dài hơn có tốc độ dài lớn hơn nếu cùng tốc độ góc.
\item (Sai) Mọi kim đồng hồ có cùng chu kì.
\item (Đúng) Đồ thị góc theo thời gian của chuyển động tròn đều là đường thẳng.
\end{itemchoice}
}
\end{ex}

% ===== Câu 149 =====
% ID: AIQ_L10C6_FULL_0147
% Mức: TH
\begin{ex}
% ID: AIQ_L10C6_FULL_0147
% Mức: TH
Xét các phát biểu thuộc dạng “Khai thác đồ thị và tình huống thực tiễn về chuyển động tròn đều” ở tình huống số 2.
\choiceTF
{\True Kim giây có chu kì $60\,\text{s}$.}
{\True Đầu kim dài hơn có tốc độ dài lớn hơn nếu cùng tốc độ góc.}
{Mọi kim đồng hồ có cùng chu kì.}
{\True Đồ thị góc theo thời gian của chuyển động tròn đều là đường thẳng.}
\loigiai{
\begin{itemchoice}
\item (Đúng) Kim giây có chu kì $60\,\text{s}$. (Vì): Phù hợp với định nghĩa hoặc hệ thức vật lí. b) Đúng: Phù hợp với định nghĩa hoặc hệ thức vật lí. c) Sai: Không phù hợp với định nghĩa hoặc hệ thức vật lí. d) Đúng: Phù hợp với định nghĩa hoặc hệ thức vật lí
\item (Đúng) Đầu kim dài hơn có tốc độ dài lớn hơn nếu cùng tốc độ góc.
\item (Sai) Mọi kim đồng hồ có cùng chu kì.
\item (Đúng) Đồ thị góc theo thời gian của chuyển động tròn đều là đường thẳng.
\end{itemchoice}
}
\end{ex}

% ===== Câu 150 =====
% ID: AIQ_L10C6_FULL_0148
% Mức: TH
\begin{ex}
% ID: AIQ_L10C6_FULL_0148
% Mức: TH
Xét các phát biểu thuộc dạng “Khai thác đồ thị và tình huống thực tiễn về chuyển động tròn đều” ở tình huống số 3.
\choiceTF
{\True Kim giây có chu kì $60\,\text{s}$.}
{\True Đầu kim dài hơn có tốc độ dài lớn hơn nếu cùng tốc độ góc.}
{Mọi kim đồng hồ có cùng chu kì.}
{\True Đồ thị góc theo thời gian của chuyển động tròn đều là đường thẳng.}
\loigiai{
\begin{itemchoice}
\item (Đúng) Kim giây có chu kì $60\,\text{s}$. (Vì): Phù hợp với định nghĩa hoặc hệ thức vật lí. b) Đúng: Phù hợp với định nghĩa hoặc hệ thức vật lí. c) Sai: Không phù hợp với định nghĩa hoặc hệ thức vật lí. d) Đúng: Phù hợp với định nghĩa hoặc hệ thức vật lí
\item (Đúng) Đầu kim dài hơn có tốc độ dài lớn hơn nếu cùng tốc độ góc.
\item (Sai) Mọi kim đồng hồ có cùng chu kì.
\item (Đúng) Đồ thị góc theo thời gian của chuyển động tròn đều là đường thẳng.
\end{itemchoice}
}
\end{ex}

% ===== Câu 151 =====
% ID: AIQ_L10C6_FULL_0149
% Mức: VD
\begin{ex}
% ID: AIQ_L10C6_FULL_0149
% Mức: VD
Xét các phát biểu thuộc dạng “Khai thác đồ thị và tình huống thực tiễn về chuyển động tròn đều” ở tình huống số 4.
\choiceTF
{\True Kim giây có chu kì $60\,\text{s}$.}
{\True Đầu kim dài hơn có tốc độ dài lớn hơn nếu cùng tốc độ góc.}
{Mọi kim đồng hồ có cùng chu kì.}
{\True Đồ thị góc theo thời gian của chuyển động tròn đều là đường thẳng.}
\loigiai{
\begin{itemchoice}
\item (Đúng) Kim giây có chu kì $60\,\text{s}$. (Vì): Phù hợp với định nghĩa hoặc hệ thức vật lí. b) Đúng: Phù hợp với định nghĩa hoặc hệ thức vật lí. c) Sai: Không phù hợp với định nghĩa hoặc hệ thức vật lí. d) Đúng: Phù hợp với định nghĩa hoặc hệ thức vật lí
\item (Đúng) Đầu kim dài hơn có tốc độ dài lớn hơn nếu cùng tốc độ góc.
\item (Sai) Mọi kim đồng hồ có cùng chu kì.
\item (Đúng) Đồ thị góc theo thời gian của chuyển động tròn đều là đường thẳng.
\end{itemchoice}
}
\end{ex}

% ===== Câu 152 =====
% ID: AIQ_L10C6_FULL_0150
% Mức: VDC
\begin{ex}
% ID: AIQ_L10C6_FULL_0150
% Mức: VDC
Xét các phát biểu thuộc dạng “Khai thác đồ thị và tình huống thực tiễn về chuyển động tròn đều” ở tình huống số 5.
\choiceTF
{\True Kim giây có chu kì $60\,\text{s}$.}
{\True Đầu kim dài hơn có tốc độ dài lớn hơn nếu cùng tốc độ góc.}
{Mọi kim đồng hồ có cùng chu kì.}
{\True Đồ thị góc theo thời gian của chuyển động tròn đều là đường thẳng.}
\loigiai{
\begin{itemchoice}
\item (Đúng) Kim giây có chu kì $60\,\text{s}$. (Vì): Phù hợp với định nghĩa hoặc hệ thức vật lí. b) Đúng: Phù hợp với định nghĩa hoặc hệ thức vật lí. c) Sai: Không phù hợp với định nghĩa hoặc hệ thức vật lí. d) Đúng: Phù hợp với định nghĩa hoặc hệ thức vật lí
\item (Đúng) Đầu kim dài hơn có tốc độ dài lớn hơn nếu cùng tốc độ góc.
\item (Sai) Mọi kim đồng hồ có cùng chu kì.
\item (Đúng) Đồ thị góc theo thời gian của chuyển động tròn đều là đường thẳng.
\end{itemchoice}
}
\end{ex}

% ===== Câu 153 =====
% ID: AIQ_L10C6_FULL_0151
% Mức: TH
\begin{ex}
% ID: AIQ_L10C6_FULL_0151
% Mức: TH
Đầu kim giây dài $0,2$ m quay đều. Tính tốc độ dài của đầu kim.
\shortans{0,02}
\loigiai{
Kim giây có $T=60$ s, nên $v=2\pi r/T=2\pi\cdot0,2/60$. Suy ra kết quả $0,02\ \mathrm{m/s}$. Đáp án trả lời ngắn không ghi đơn vị.
}
\end{ex}

% ===== Câu 154 =====
% ID: AIQ_L10C6_FULL_0152
% Mức: TH
\begin{ex}
% ID: AIQ_L10C6_FULL_0152
% Mức: TH
Đầu kim giây dài $0,21$ m quay đều. Tính tốc độ dài của đầu kim.
\shortans{0,02}
\loigiai{
Kim giây có $T=60$ s, nên $v=2\pi r/T=2\pi\cdot0,21/60$. Suy ra kết quả $0,02\ \mathrm{m/s}$. Đáp án trả lời ngắn không ghi đơn vị.
}
\end{ex}

% ===== Câu 155 =====
% ID: AIQ_L10C6_FULL_0153
% Mức: VD
\begin{ex}
% ID: AIQ_L10C6_FULL_0153
% Mức: VD
Đầu kim giây dài $0,22$ m quay đều. Tính tốc độ dài của đầu kim.
\shortans{0,02}
\loigiai{
Kim giây có $T=60$ s, nên $v=2\pi r/T=2\pi\cdot0,22/60$. Suy ra kết quả $0,02\ \mathrm{m/s}$. Đáp án trả lời ngắn không ghi đơn vị.
}
\end{ex}

% ===== Câu 156 =====
% ID: AIQ_L10C6_FULL_0154
% Mức: VD
\begin{ex}
% ID: AIQ_L10C6_FULL_0154
% Mức: VD
Đầu kim giây dài $0,23$ m quay đều. Tính tốc độ dài của đầu kim.
\shortans{0,02}
\loigiai{
Kim giây có $T=60$ s, nên $v=2\pi r/T=2\pi\cdot0,23/60$. Suy ra kết quả $0,02\ \mathrm{m/s}$. Đáp án trả lời ngắn không ghi đơn vị.
}
\end{ex}

% ===== Câu 157 =====
% ID: AIQ_L10C6_FULL_0155
% Mức: VDC
\begin{ex}
% ID: AIQ_L10C6_FULL_0155
% Mức: VDC
Đầu kim giây dài $0,24$ m quay đều. Tính tốc độ dài của đầu kim.
\shortans{0,03}
\loigiai{
Kim giây có $T=60$ s, nên $v=2\pi r/T=2\pi\cdot0,24/60$. Suy ra kết quả $0,03\ \mathrm{m/s}$. Đáp án trả lời ngắn không ghi đơn vị.
}
\end{ex}

% ===== Câu 158 =====
% ID: AIQ_L10C6_FULL_0156
% Mức: VDC
\begin{ex}
% ID: AIQ_L10C6_FULL_0156
% Mức: VDC
Đầu kim giây dài $0,25$ m quay đều. Tính tốc độ dài của đầu kim.
\shortans{0,03}
\loigiai{
Kim giây có $T=60$ s, nên $v=2\pi r/T=2\pi\cdot0,25/60$. Suy ra kết quả $0,03\ \mathrm{m/s}$. Đáp án trả lời ngắn không ghi đơn vị.
}
\end{ex}

% ===== Câu 159 =====
% ID: AIQ_L10C6_FULL_0157
% Mức: TH
\begin{bt}
% ID: AIQ_L10C6_FULL_0157
% Mức: TH
Trình bày cách giải: Đầu kim giây dài $0,26$ m quay đều. Tính tốc độ dài của đầu kim.
\loigiai{
Kim giây có $T=60$ s, nên $v=2\pi r/T=2\pi\cdot0,26/60$. Suy ra kết quả $0,03\ \mathrm{m/s}$.
}
\end{bt}

% ===== Câu 160 =====
% ID: AIQ_L10C6_FULL_0158
% Mức: TH
\begin{bt}
% ID: AIQ_L10C6_FULL_0158
% Mức: TH
Trình bày cách giải: Đầu kim giây dài $0,27$ m quay đều. Tính tốc độ dài của đầu kim.
\loigiai{
Kim giây có $T=60$ s, nên $v=2\pi r/T=2\pi\cdot0,27/60$. Suy ra kết quả $0,03\ \mathrm{m/s}$.
}
\end{bt}

% ===== Câu 161 =====
% ID: AIQ_L10C6_FULL_0159
% Mức: VD
\begin{bt}
% ID: AIQ_L10C6_FULL_0159
% Mức: VD
Trình bày cách giải: Đầu kim giây dài $0,28$ m quay đều. Tính tốc độ dài của đầu kim.
\loigiai{
Kim giây có $T=60$ s, nên $v=2\pi r/T=2\pi\cdot0,28/60$. Suy ra kết quả $0,03\ \mathrm{m/s}$.
}
\end{bt}

% ===== Câu 162 =====
% ID: AIQ_L10C6_FULL_0160
% Mức: VD
\begin{bt}
% ID: AIQ_L10C6_FULL_0160
% Mức: VD
Trình bày cách giải: Đầu kim giây dài $0,29$ m quay đều. Tính tốc độ dài của đầu kim.
\loigiai{
Kim giây có $T=60$ s, nên $v=2\pi r/T=2\pi\cdot0,29/60$. Suy ra kết quả $0,03\ \mathrm{m/s}$.
}
\end{bt}

% ===== Câu 163 =====
% ID: AIQ_L10C6_FULL_0161
% Mức: VDC
\begin{bt}
% ID: AIQ_L10C6_FULL_0161
% Mức: VDC
Trình bày cách giải: Đầu kim giây dài $0,3$ m quay đều. Tính tốc độ dài của đầu kim.
\loigiai{
Kim giây có $T=60$ s, nên $v=2\pi r/T=2\pi\cdot0,3/60$. Suy ra kết quả $0,03\ \mathrm{m/s}$.
}
\end{bt}

% ===== Câu 164 =====
% ID: AIQ_L10C6_FULL_0162
% Mức: VDC
\begin{bt}
% ID: AIQ_L10C6_FULL_0162
% Mức: VDC
Trình bày cách giải: Đầu kim giây dài $0,31$ m quay đều. Tính tốc độ dài của đầu kim.
\loigiai{
Kim giây có $T=60$ s, nên $v=2\pi r/T=2\pi\cdot0,31/60$. Suy ra kết quả $0,03\ \mathrm{m/s}$.
}
\end{bt}
